
\documentclass{article}
\usepackage{colortbl}
\usepackage{makecell}
\usepackage{multirow}
\usepackage{supertabular}

\begin{document}

\newcounter{utterance}

\centering \large Interaction Transcript for game `clean\_up', experiment `1\_medium\_7obj\_ar', episode 3 with nemotron{-}3{-}ultra.
\vspace{24pt}

{ \footnotesize  \setcounter{utterance}{1}
\setlength{\tabcolsep}{0pt}
\begin{supertabular}{c@{$\;$}|p{.15\linewidth}@{}p{.15\linewidth}p{.15\linewidth}p{.15\linewidth}p{.15\linewidth}p{.15\linewidth}}
    \# & $\;$A & \multicolumn{4}{c}{Game Master} & $\;\:$B\\
    \hline

    \theutterance \stepcounter{utterance}  
    & & \multicolumn{4}{p{0.6\linewidth}}{
        \cellcolor[rgb]{0.9,0.9,0.9}{
            \makecell[{{p{\linewidth}}}]{
                \texttt{\tiny{[P1$\langle$GM]}}
                \texttt{أنا مدير لعبتك، وأنت تلعب لعبة تعاونية مع الشبكة التالية كلوحة لعب:} \\
\\ 
\texttt{```} \\
\texttt{ 1234567} \\
\texttt{╔═╤═════╗} \\
\texttt{║◌│◌A◌◌◌║ 1} \\
\texttt{╟─┤◌T◌◌◌║ 2} \\
\texttt{║O└─┐◌◌◌║ 3} \\
\texttt{║◌◌◌├─┐◌║ 4} \\
\texttt{║◌◌S│◌│◌║ 5} \\
\texttt{╟─┬─┴─┼─╢ 6} \\
\texttt{║P│◌DM│◌║ 7} \\
\texttt{╚═╧═══╧═╝} \\
\\ 
\texttt{```} \\
\\ 
\texttt{* تعرض الحافة العلوية إحداثيات x التي تتزايد نحو اليمين، والحافة اليمنى إحداثيات y التي تتزايد نحو الأسفل.} \\
\texttt{* الكائنات التالية موضوعة عشوائياً على شبكتك: 'P', 'O', 'T', 'S', 'D', 'A', 'M'.} \\
\\ 
\\ 
\texttt{يرى اللاعب الآخر نسخة مختلفة من لوحة اللعب، حيث توضع الكائنات في مواقع عشوائية مختلفة. لا يمكنك رؤية لوحة اللاعب الآخر، ولا يمكنه رؤية لوحتك.} \\
\\ 
\\ 
\texttt{**الهدف:**} \\
\\ 
\texttt{يجب على كلا اللاعبين تحريك الكائنات على خلفية كل منهما بحيث تنتهي الكائنات المتطابقة عند الإحداثيات نفسها. عليك التواصل مع اللاعب الآخر للاتفاق على إعداد هدف مشترك.} \\
\\ 
\\ 
\texttt{**القواعد:**} \\
\\ 
\texttt{* في كل دور، يمكنك إرسال أمر واحد فقط من الأمرين التاليين:} \\
\texttt{ 1. `قل: <رسالة>`: لإرسال رسالة (كل شيء حتى فاصل السطر التالي) إلى اللاعب الآخر. سأقوم بتمريرها إلى شريكك.} \\
\texttt{2. `حرّك: <كائن>, (<X>, <Y>)`: لتحريك كائن إلى موضع جديد، حيث `<X>` هو العمود و `<Y>` هو الصف. سأخبرك بما إذا كانت حركتك ناجحة أم لا.} \\
\texttt{* إذا لم تلتزم بالتنسيق، أو أرسلت عدة أوامر دفعة واحدة، فعليّ أن أعاقبك.} \\
\texttt{* إذا تجاوز كلا اللاعبين أكثر من 16 عقوبة، فستخسران اللعبة معاً.} \\
\texttt{* من الضروري أن تتواصل مع اللاعب الآخر بشأن حالة هدفك! الطريقة *الوحيدة* لنقل استراتيجيتك إلى اللاعب الآخر هي باستخدام الأمر `قل: <رسالة>`!} \\
\\ 
\\ 
\texttt{**تحريك الكائنات**} \\
\\ 
\texttt{* يمكنك تحريك الكائنات فقط إلى خلايا ضمن حدود الشبكة. يجب أن تكون الخلية المستهدفة فارغة، أي يجب أن تحتوي فقط على الرمز '◌'.} \\
\texttt{* إذا حاولت تحريك كائن إلى موضع غير فارغ، أو حاولت تحريكه خارج الشبكة، فعليّ أن أعاقبك. ستحصل على محاولة أخرى.} \\
\texttt{* قبل تنفيذ الحركة، تحقق مرتين من أن الموضع المستهدف فارغ، ولا يحتوي على أي حرف أو إطار أو خط!} \\
\\ 
\\ 
\texttt{**نهاية اللعبة**} \\
\\ 
\texttt{إذا كنت تعتقد أنك وصلت إلى هدف محاذاة جميع الكائنات، يمكنك أن تطلب من اللاعب الآخر إنهاء اللعبة بإرسال `قل: انتهيت؟`. إذا طلب منك اللاعب الآخر إنهاء اللعبة، ورددت بـ `قل: انتهيت!`، ستنتهي اللعبة.} \\
\\ 
\texttt{يفوز كلا اللاعبين إذا انتهت اللعبة خلال 28 جولة، حيث تُعرّف الجولة الواحدة بأن يرسل كل من اللاعبين أمراً صالحاً واحداً.} \\
\\ 
\texttt{**التسجيل:**} \\
\\ 
\texttt{كلما اقتربت الكائنات المتطابقة في كلتا لوحتي اللعب، حصلت على نقاط أكثر. العقوبات تقلل نقاطك. هل يمكنك تحطيم الرقم القياسي؟} \\
\\ 
\\ 
\texttt{يرجى إرسال رسالة إلى اللاعب الآخر لبدء اللعبة!} \\
            }
        }
    }
    & \\ \\

    \theutterance \stepcounter{utterance}  
    & \multicolumn{4}{p{0.6\linewidth}}{
        \cellcolor[rgb]{0.9,0.9,0.9}{
            \makecell[{{p{\linewidth}}}]{
                \texttt{\tiny{[P1$\rangle$GM]}}
                \texttt{سأبدأ بالتواصل مع اللاعب الآخر لفهم مواقع الكائنات لديهم والاتفاق على استراتيجية.} \\
\\ 
\texttt{قل: مرحباً! أنا لاعب 1. لبدء التنسيق، يرجى إخباري بإحداثيات جميع الكائنات على لوحتك (P, O, T, S, D, A, M). على لوحتي، المواقع هي:} \\
\texttt{P: (1,7)} \\
\texttt{O: (1,3)} \\
\texttt{T: (3,2)} \\
\texttt{S: (3,5)} \\
\texttt{D: (5,7)} \\
\texttt{A: (3,1)} \\
\texttt{M: (6,7)} \\
\\ 
\texttt{ما هي مواقعك؟} \\
            }
        }
    }
    & & \\ \\

    \theutterance \stepcounter{utterance}  
    & & & \multicolumn{2}{p{0.3\linewidth}}{
        \cellcolor[rgb]{0.9,0.9,0.9}{
            \makecell[{{p{\linewidth}}}]{
                \texttt{\tiny{[GM$|$GM]}}
                \texttt{سأبدأ بالتواصل مع اللاعب الآخر لفهم مواقع الكائنات لديهم والاتفاق على استراتيجية.} \\
\\ 
\texttt{قل: مرحباً! أنا لاعب 1. لبدء التنسيق، يرجى إخباري بإحداثيات جميع الكائنات على لوحتك (P, O, T, S, D, A, M). على لوحتي، المواقع هي:} \\
\texttt{P: (1,7)} \\
\texttt{O: (1,3)} \\
\texttt{T: (3,2)} \\
\texttt{S: (3,5)} \\
\texttt{D: (5,7)} \\
\texttt{A: (3,1)} \\
\texttt{M: (6,7)} \\
\\ 
\texttt{ما هي مواقعك؟} \\
            }
        }
    }
    & & \\ \\

    \theutterance \stepcounter{utterance}  
    & & & \multicolumn{2}{p{0.3\linewidth}}{
        \cellcolor[rgb]{0.9,0.9,0.9}{
            \makecell[{{p{\linewidth}}}]{
                \texttt{\tiny{[GM$|$GM]}}
                \texttt{Invalid format: head and tail are not empty} \\
\texttt{head: 'سأبدأ بالتواصل مع اللاعب الآخر لفهم مواقع الكائنات لديهم والاتفاق على استراتيجية.} \\
\\ 
\texttt{'} \\
\texttt{tail: '} \\
\texttt{P: (1,7)} \\
\texttt{O: (1,3)} \\
\texttt{T: (3,2)} \\
\texttt{S: (3,5)} \\
\texttt{D: (5,7)} \\
\texttt{A: (3,1)} \\
\texttt{M: (6,7)} \\
\\ 
\texttt{ما هي مواقعك؟'} \\
            }
        }
    }
    & & \\ \\

    \theutterance \stepcounter{utterance}  
    & & \multicolumn{4}{p{0.6\linewidth}}{
        \cellcolor[rgb]{0.9,0.9,0.9}{
            \makecell[{{p{\linewidth}}}]{
                \texttt{\tiny{[P1$\langle$GM]}}
                \texttt{عقوبة: يجب ألا تحتوي رسالتك على أي شيء قبل الأمر أو بعده!} \\
\\ 
\texttt{تأكد من أن ردك يحتوي فقط على `قل: <رسالة>` أو `حرّك: <كائن>, (<X>, <Y>)`، ولا شيء غير ذلك!} \\
\\ 
\texttt{لقد جمعتما معاً 1 من أصل 16 عقوبة.} \\
\texttt{يرجى المحاولة مرة أخرى!} \\
            }
        }
    }
    & \\ \\

    \theutterance \stepcounter{utterance}  
    & \multicolumn{4}{p{0.6\linewidth}}{
        \cellcolor[rgb]{0.9,0.9,0.9}{
            \makecell[{{p{\linewidth}}}]{
                \texttt{\tiny{[P1$\rangle$GM]}}
                \texttt{قل: مرحباً! أنا لاعب 1. لبدء التنسيق، يرجى إخباري بإحداثيات جميع الكائنات على لوحتك (P, O, T, S, D, A, M). على لوحتي، المواقع هي:} \\
\texttt{P: (1,7)} \\
\texttt{O: (1,3)} \\
\texttt{T: (3,2)} \\
\texttt{S: (3,5)} \\
\texttt{D: (5,7)} \\
\texttt{A: (3,1)} \\
\texttt{M: (6,7)} \\
\\ 
\texttt{ما هي مواقعك؟} \\
            }
        }
    }
    & & \\ \\

    \theutterance \stepcounter{utterance}  
    & & & \multicolumn{2}{p{0.3\linewidth}}{
        \cellcolor[rgb]{0.9,0.9,0.9}{
            \makecell[{{p{\linewidth}}}]{
                \texttt{\tiny{[GM$|$GM]}}
                \texttt{قل: مرحباً! أنا لاعب 1. لبدء التنسيق، يرجى إخباري بإحداثيات جميع الكائنات على لوحتك (P, O, T, S, D, A, M). على لوحتي، المواقع هي:} \\
\texttt{P: (1,7)} \\
\texttt{O: (1,3)} \\
\texttt{T: (3,2)} \\
\texttt{S: (3,5)} \\
\texttt{D: (5,7)} \\
\texttt{A: (3,1)} \\
\texttt{M: (6,7)} \\
\\ 
\texttt{ما هي مواقعك؟} \\
            }
        }
    }
    & & \\ \\

    \theutterance \stepcounter{utterance}  
    & & & \multicolumn{2}{p{0.3\linewidth}}{
        \cellcolor[rgb]{0.9,0.9,0.9}{
            \makecell[{{p{\linewidth}}}]{
                \texttt{\tiny{[GM$|$GM]}}
                \texttt{Invalid format: tail is not empty: '} \\
\texttt{P: (1,7)} \\
\texttt{O: (1,3)} \\
\texttt{T: (3,2)} \\
\texttt{S: (3,5)} \\
\texttt{D: (5,7)} \\
\texttt{A: (3,1)} \\
\texttt{M: (6,7)} \\
\\ 
\texttt{ما هي مواقعك؟} \\
            }
        }
    }
    & & \\ \\

    \theutterance \stepcounter{utterance}  
    & & \multicolumn{4}{p{0.6\linewidth}}{
        \cellcolor[rgb]{0.9,0.9,0.9}{
            \makecell[{{p{\linewidth}}}]{
                \texttt{\tiny{[P1$\langle$GM]}}
                \texttt{عقوبة: يجب ألا تحتوي رسالتك على أي شيء بعد الأمر!} \\
\\ 
\texttt{تأكد من أن ردك يحتوي فقط على `قل: <رسالة>` أو `حرّك: <كائن>, (<X>, <Y>)`، ولا شيء غير ذلك!} \\
\\ 
\texttt{لقد جمعتما معاً 2 من أصل 16 عقوبة.} \\
\texttt{يرجى المحاولة مرة أخرى!} \\
            }
        }
    }
    & \\ \\

    \theutterance \stepcounter{utterance}  
    & \multicolumn{4}{p{0.6\linewidth}}{
        \cellcolor[rgb]{0.9,0.9,0.9}{
            \makecell[{{p{\linewidth}}}]{
                \texttt{\tiny{[P1$\rangle$GM]}}
                \texttt{قل: مرحباً! أنا لاعب 1. مواقع كائناتي: P(1,7) O(1,3) T(3,2) S(3,5) D(5,7) A(3,1) M(6,7). ما هي مواقعك؟} \\
            }
        }
    }
    & & \\ \\

    \theutterance \stepcounter{utterance}  
    & & & \multicolumn{2}{p{0.3\linewidth}}{
        \cellcolor[rgb]{0.9,0.9,0.9}{
            \makecell[{{p{\linewidth}}}]{
                \texttt{\tiny{[GM$|$GM]}}
                \texttt{قل: مرحباً! أنا لاعب 1. مواقع كائناتي: P(1,7) O(1,3) T(3,2) S(3,5) D(5,7) A(3,1) M(6,7). ما هي مواقعك؟} \\
            }
        }
    }
    & & \\ \\

    \theutterance \stepcounter{utterance}  
    & & \multicolumn{4}{p{0.6\linewidth}}{
        \cellcolor[rgb]{0.9,0.9,0.9}{
            \makecell[{{p{\linewidth}}}]{
                \texttt{\tiny{[GM$\rangle$P2]}}
                \texttt{أنا مدير لعبتك، وأنت تلعب لعبة تعاونية مع الشبكة التالية كلوحة لعب:} \\
\\ 
\texttt{```} \\
\texttt{ 1234567} \\
\texttt{╔═╤═════╗} \\
\texttt{║P│S◌D◌◌║ 1} \\
\texttt{╟─┤◌◌◌◌◌║ 2} \\
\texttt{║◌└─┐◌A◌║ 3} \\
\texttt{║◌◌M├─┐◌║ 4} \\
\texttt{║◌T◌│◌│O║ 5} \\
\texttt{╟─┬─┴─┼─╢ 6} \\
\texttt{║◌│◌◌◌│◌║ 7} \\
\texttt{╚═╧═══╧═╝} \\
\\ 
\texttt{```} \\
\\ 
\texttt{* تعرض الحافة العلوية إحداثيات x التي تتزايد نحو اليمين، والحافة اليمنى إحداثيات y التي تتزايد نحو الأسفل.} \\
\texttt{* الكائنات التالية موضوعة عشوائياً على شبكتك: 'P', 'O', 'T', 'S', 'D', 'A', 'M'.} \\
\\ 
\\ 
\texttt{يرى اللاعب الآخر نسخة مختلفة من لوحة اللعب، حيث توضع الكائنات في مواقع عشوائية مختلفة. لا يمكنك رؤية لوحة اللاعب الآخر، ولا يمكنه رؤية لوحتك.} \\
\\ 
\\ 
\texttt{**الهدف:**} \\
\\ 
\texttt{يجب على كلا اللاعبين تحريك الكائنات على خلفية كل منهما بحيث تنتهي الكائنات المتطابقة عند الإحداثيات نفسها. عليك التواصل مع اللاعب الآخر للاتفاق على إعداد هدف مشترك.} \\
\\ 
\\ 
\texttt{**القواعد:**} \\
\\ 
\texttt{* في كل دور، يمكنك إرسال أمر واحد فقط من الأمرين التاليين:} \\
\texttt{ 1. `قل: <رسالة>`: لإرسال رسالة (كل شيء حتى فاصل السطر التالي) إلى اللاعب الآخر. سأقوم بتمريرها إلى شريكك.} \\
\texttt{2. `حرّك: <كائن>, (<X>, <Y>)`: لتحريك كائن إلى موضع جديد، حيث `<X>` هو العمود و `<Y>` هو الصف. سأخبرك بما إذا كانت حركتك ناجحة أم لا.} \\
\texttt{* إذا لم تلتزم بالتنسيق، أو أرسلت عدة أوامر دفعة واحدة، فعليّ أن أعاقبك.} \\
\texttt{* إذا تجاوز كلا اللاعبين أكثر من 16 عقوبة، فستخسران اللعبة معاً.} \\
\texttt{* من الضروري أن تتواصل مع اللاعب الآخر بشأن حالة هدفك! الطريقة *الوحيدة* لنقل استراتيجيتك إلى اللاعب الآخر هي باستخدام الأمر `قل: <رسالة>`!} \\
\\ 
\\ 
\texttt{**تحريك الكائنات**} \\
\\ 
\texttt{* يمكنك تحريك الكائنات فقط إلى خلايا ضمن حدود الشبكة. يجب أن تكون الخلية المستهدفة فارغة، أي يجب أن تحتوي فقط على الرمز '◌'.} \\
\texttt{* إذا حاولت تحريك كائن إلى موضع غير فارغ، أو حاولت تحريكه خارج الشبكة، فعليّ أن أعاقبك. ستحصل على محاولة أخرى.} \\
\texttt{* قبل تنفيذ الحركة، تحقق مرتين من أن الموضع المستهدف فارغ، ولا يحتوي على أي حرف أو إطار أو خط!} \\
\\ 
\\ 
\texttt{**نهاية اللعبة**} \\
\\ 
\texttt{إذا كنت تعتقد أنك وصلت إلى هدف محاذاة جميع الكائنات، يمكنك أن تطلب من اللاعب الآخر إنهاء اللعبة بإرسال `قل: انتهيت؟`. إذا طلب منك اللاعب الآخر إنهاء اللعبة، ورددت بـ `قل: انتهيت!`، ستنتهي اللعبة.} \\
\\ 
\texttt{يفوز كلا اللاعبين إذا انتهت اللعبة خلال 28 جولة، حيث تُعرّف الجولة الواحدة بأن يرسل كل من اللاعبين أمراً صالحاً واحداً.} \\
\\ 
\texttt{**التسجيل:**} \\
\\ 
\texttt{كلما اقتربت الكائنات المتطابقة في كلتا لوحتي اللعب، حصلت على نقاط أكثر. العقوبات تقلل نقاطك. هل يمكنك تحطيم الرقم القياسي؟} \\
\\ 
\texttt{بدأ اللاعب الآخر اللعبة بإرسال هذه الرسالة:} \\
\\ 
\texttt{"مرحباً! أنا لاعب 1. مواقع كائناتي: P(1,7) O(1,3) T(3,2) S(3,5) D(5,7) A(3,1) M(6,7). ما هي مواقعك؟"} \\
\\ 
\texttt{ما هو أمرك الأول؟} \\
            }
        }
    }
    & \\ \\

    \theutterance \stepcounter{utterance}  
    & & & \multicolumn{4}{p{0.6\linewidth}}{
        \cellcolor[rgb]{0.9,0.9,0.9}{
            \makecell[{{p{\linewidth}}}]{
                \texttt{\tiny{[GM$\langle$P2]}}
                \texttt{قول: مرحباً! أنا لاعب 2. مواقع كائناتي: P(1,1) S(2,1) D(4,1) A(5,3) M(3,4) T(2,5) O(7,5).} \\
            }
        }
    }
     \\ \\

    \theutterance \stepcounter{utterance}  
    & & & \multicolumn{2}{p{0.3\linewidth}}{
        \cellcolor[rgb]{0.9,0.9,0.9}{
            \makecell[{{p{\linewidth}}}]{
                \texttt{\tiny{[GM$|$GM]}}
                \texttt{قول: مرحباً! أنا لاعب 2. مواقع كائناتي: P(1,1) S(2,1) D(4,1) A(5,3) M(3,4) T(2,5) O(7,5).} \\
            }
        }
    }
    & & \\ \\

    \theutterance \stepcounter{utterance}  
    & & & \multicolumn{2}{p{0.3\linewidth}}{
        \cellcolor[rgb]{0.9,0.9,0.9}{
            \makecell[{{p{\linewidth}}}]{
                \texttt{\tiny{[GM$|$GM]}}
                \texttt{Invalid response format} \\
            }
        }
    }
    & & \\ \\

    \theutterance \stepcounter{utterance}  
    & & \multicolumn{4}{p{0.6\linewidth}}{
        \cellcolor[rgb]{0.9,0.9,0.9}{
            \makecell[{{p{\linewidth}}}]{
                \texttt{\tiny{[GM$\rangle$P2]}}
                \texttt{عقوبة: رسالتك ليست بالتنسيق المتوقع!} \\
\\ 
\texttt{تأكد من أن ردك يحتوي فقط على `قل: <رسالة>` أو `حرّك: <كائن>, (<X>, <Y>)`، ولا شيء غير ذلك!} \\
\\ 
\texttt{لقد جمعتما معاً 3 من أصل 16 عقوبة.} \\
\texttt{يرجى المحاولة مرة أخرى!} \\
            }
        }
    }
    & \\ \\

    \theutterance \stepcounter{utterance}  
    & & & \multicolumn{4}{p{0.6\linewidth}}{
        \cellcolor[rgb]{0.9,0.9,0.9}{
            \makecell[{{p{\linewidth}}}]{
                \texttt{\tiny{[GM$\langle$P2]}}
                \texttt{قول: مرحباً! أنا لاعب 2. مواقع كائناتي: P(1,1) S(2,1) D(4,1) A(5,3) M(3,4) T(2,5) O(7,5)} \\
            }
        }
    }
     \\ \\

    \theutterance \stepcounter{utterance}  
    & & & \multicolumn{2}{p{0.3\linewidth}}{
        \cellcolor[rgb]{0.9,0.9,0.9}{
            \makecell[{{p{\linewidth}}}]{
                \texttt{\tiny{[GM$|$GM]}}
                \texttt{قول: مرحباً! أنا لاعب 2. مواقع كائناتي: P(1,1) S(2,1) D(4,1) A(5,3) M(3,4) T(2,5) O(7,5)} \\
            }
        }
    }
    & & \\ \\

    \theutterance \stepcounter{utterance}  
    & & & \multicolumn{2}{p{0.3\linewidth}}{
        \cellcolor[rgb]{0.9,0.9,0.9}{
            \makecell[{{p{\linewidth}}}]{
                \texttt{\tiny{[GM$|$GM]}}
                \texttt{Invalid response format} \\
            }
        }
    }
    & & \\ \\

    \theutterance \stepcounter{utterance}  
    & & \multicolumn{4}{p{0.6\linewidth}}{
        \cellcolor[rgb]{0.9,0.9,0.9}{
            \makecell[{{p{\linewidth}}}]{
                \texttt{\tiny{[GM$\rangle$P2]}}
                \texttt{عقوبة: رسالتك ليست بالتنسيق المتوقع!} \\
\\ 
\texttt{تأكد من أن ردك يحتوي فقط على `قل: <رسالة>` أو `حرّك: <كائن>, (<X>, <Y>)`، ولا شيء غير ذلك!} \\
\\ 
\texttt{لقد جمعتما معاً 4 من أصل 16 عقوبة.} \\
\texttt{يرجى المحاولة مرة أخرى!} \\
            }
        }
    }
    & \\ \\

    \theutterance \stepcounter{utterance}  
    & & & \multicolumn{4}{p{0.6\linewidth}}{
        \cellcolor[rgb]{0.9,0.9,0.9}{
            \makecell[{{p{\linewidth}}}]{
                \texttt{\tiny{[GM$\langle$P2]}}
                \texttt{قل: مرحباً! أنا لاعب 2. مواقع كائناتي: P(1,1) S(2,1) D(4,1) A(5,3) M(3,4) T(2,5) O(7,5)} \\
            }
        }
    }
     \\ \\

    \theutterance \stepcounter{utterance}  
    & & & \multicolumn{2}{p{0.3\linewidth}}{
        \cellcolor[rgb]{0.9,0.9,0.9}{
            \makecell[{{p{\linewidth}}}]{
                \texttt{\tiny{[GM$|$GM]}}
                \texttt{قل: مرحباً! أنا لاعب 2. مواقع كائناتي: P(1,1) S(2,1) D(4,1) A(5,3) M(3,4) T(2,5) O(7,5)} \\
            }
        }
    }
    & & \\ \\

    \theutterance \stepcounter{utterance}  
    & & \multicolumn{4}{p{0.6\linewidth}}{
        \cellcolor[rgb]{0.9,0.9,0.9}{
            \makecell[{{p{\linewidth}}}]{
                \texttt{\tiny{[P1$\langle$GM]}}
                \texttt{تم تمرير رسالتك إلى اللاعب الآخر.} \\
\\ 
\\ 
\texttt{أنت تلعب حالياً الجولة 1 من أصل 28 كحد أقصى.} \\
\texttt{لقد جمعتما معاً 4 من أصل 16 عقوبة.} \\
\texttt{أرسل اللاعب الآخر هذه الرسالة:} \\
\\ 
\\ 
\texttt{"مرحباً! أنا لاعب 2. مواقع كائناتي: P(1,1) S(2,1) D(4,1) A(5,3) M(3,4) T(2,5) O(7,5)"} \\
\\ 
\\ 
\texttt{ما هو أمرك التالي؟} \\
            }
        }
    }
    & \\ \\

    \theutterance \stepcounter{utterance}  
    & \multicolumn{4}{p{0.6\linewidth}}{
        \cellcolor[rgb]{0.9,0.9,0.9}{
            \makecell[{{p{\linewidth}}}]{
                \texttt{\tiny{[P1$\rangle$GM]}}
                \texttt{قل: شكراً! مواقعك: P(1,1) S(2,1) D(4,1) A(5,3) M(3,4) T(2,5) O(7,5). } \\
\\ 
\texttt{لدينا 7 خلايا فارغة فقط على كل لوحة. أقترح أن نختار 7 إحداثيات هدف ونحرك كائناتنا إليها. الخلايا الفارغة على لوحتي: (2,1) (4,1) (5,1) (6,1) (7,1) (2,3) (4,3) (5,3) (6,3) (7,3) (1,4) (2,4) (4,4) (5,4) (6,4) (7,4) (1,5) (2,5) (4,5) (5,5) (6,5) (7,5) (1,6) (2,6) (3,6) (4,6) (7,6) (2,7) (3,7) (4,7).} \\
\\ 
\texttt{ماذا تقترح كإحداثيات هدف؟} \\
            }
        }
    }
    & & \\ \\

    \theutterance \stepcounter{utterance}  
    & & & \multicolumn{2}{p{0.3\linewidth}}{
        \cellcolor[rgb]{0.9,0.9,0.9}{
            \makecell[{{p{\linewidth}}}]{
                \texttt{\tiny{[GM$|$GM]}}
                \texttt{قل: شكراً! مواقعك: P(1,1) S(2,1) D(4,1) A(5,3) M(3,4) T(2,5) O(7,5). } \\
\\ 
\texttt{لدينا 7 خلايا فارغة فقط على كل لوحة. أقترح أن نختار 7 إحداثيات هدف ونحرك كائناتنا إليها. الخلايا الفارغة على لوحتي: (2,1) (4,1) (5,1) (6,1) (7,1) (2,3) (4,3) (5,3) (6,3) (7,3) (1,4) (2,4) (4,4) (5,4) (6,4) (7,4) (1,5) (2,5) (4,5) (5,5) (6,5) (7,5) (1,6) (2,6) (3,6) (4,6) (7,6) (2,7) (3,7) (4,7).} \\
\\ 
\texttt{ماذا تقترح كإحداثيات هدف؟} \\
            }
        }
    }
    & & \\ \\

    \theutterance \stepcounter{utterance}  
    & & & \multicolumn{2}{p{0.3\linewidth}}{
        \cellcolor[rgb]{0.9,0.9,0.9}{
            \makecell[{{p{\linewidth}}}]{
                \texttt{\tiny{[GM$|$GM]}}
                \texttt{Invalid format: tail is not empty: '} \\
\\ 
\texttt{لدينا 7 خلايا فارغة فقط على كل لوحة. أقترح أن نختار 7 إحداثيات هدف ونحرك كائناتنا إليها. الخلايا الفارغة على لوحتي: (2,1) (4,1) (5,1) (6,1) (7,1) (2,3) (4,3) (5,3) (6,3) (7,3) (1,4) (2,4) (4,4) (5,4) (6,4) (7,4) (1,5) (2,5) (4,5) (5,5) (6,5) (7,5) (1,6) (2,6) (3,6) (4,6) (7,6) (2,7) (3,7) (4,7).} \\
\\ 
\texttt{ماذا تقترح كإحداثيات هدف؟} \\
            }
        }
    }
    & & \\ \\

    \theutterance \stepcounter{utterance}  
    & & \multicolumn{4}{p{0.6\linewidth}}{
        \cellcolor[rgb]{0.9,0.9,0.9}{
            \makecell[{{p{\linewidth}}}]{
                \texttt{\tiny{[P1$\langle$GM]}}
                \texttt{عقوبة: يجب ألا تحتوي رسالتك على أي شيء بعد الأمر!} \\
\\ 
\texttt{تأكد من أن ردك يحتوي فقط على `قل: <رسالة>` أو `حرّك: <كائن>, (<X>, <Y>)`، ولا شيء غير ذلك!} \\
\\ 
\texttt{لقد جمعتما معاً 5 من أصل 16 عقوبة.} \\
\texttt{يرجى المحاولة مرة أخرى!} \\
            }
        }
    }
    & \\ \\

    \theutterance \stepcounter{utterance}  
    & \multicolumn{4}{p{0.6\linewidth}}{
        \cellcolor[rgb]{0.9,0.9,0.9}{
            \makecell[{{p{\linewidth}}}]{
                \texttt{\tiny{[P1$\rangle$GM]}}
                \texttt{قل: شكراً! مواقعك: P(1,1) S(2,1) D(4,1) A(5,3) M(3,4) T(2,5) O(7,5). أقترح استخدام الصف 1 كصف هدف لأنه فارغ تماماً على لوحتي. إحداثيات هدف مقترحة: P(1,1) S(2,1) D(3,1) A(4,1) M(5,1) T(6,1) O(7,1). هل توافق؟} \\
            }
        }
    }
    & & \\ \\

    \theutterance \stepcounter{utterance}  
    & & & \multicolumn{2}{p{0.3\linewidth}}{
        \cellcolor[rgb]{0.9,0.9,0.9}{
            \makecell[{{p{\linewidth}}}]{
                \texttt{\tiny{[GM$|$GM]}}
                \texttt{قل: شكراً! مواقعك: P(1,1) S(2,1) D(4,1) A(5,3) M(3,4) T(2,5) O(7,5). أقترح استخدام الصف 1 كصف هدف لأنه فارغ تماماً على لوحتي. إحداثيات هدف مقترحة: P(1,1) S(2,1) D(3,1) A(4,1) M(5,1) T(6,1) O(7,1). هل توافق؟} \\
            }
        }
    }
    & & \\ \\

    \theutterance \stepcounter{utterance}  
    & & \multicolumn{4}{p{0.6\linewidth}}{
        \cellcolor[rgb]{0.9,0.9,0.9}{
            \makecell[{{p{\linewidth}}}]{
                \texttt{\tiny{[GM$\rangle$P2]}}
                \texttt{تم تمرير رسالتك إلى اللاعب الآخر.} \\
\\ 
\\ 
\texttt{أنت تلعب حالياً الجولة 2 من أصل 28 كحد أقصى.} \\
\texttt{لقد جمعتما معاً 5 من أصل 16 عقوبة.} \\
\texttt{أرسل اللاعب الآخر هذه الرسالة:} \\
\\ 
\\ 
\texttt{"شكراً! مواقعك: P(1,1) S(2,1) D(4,1) A(5,3) M(3,4) T(2,5) O(7,5). أقترح استخدام الصف 1 كصف هدف لأنه فارغ تماماً على لوحتي. إحداثيات هدف مقترحة: P(1,1) S(2,1) D(3,1) A(4,1) M(5,1) T(6,1) O(7,1). هل توافق؟"} \\
\\ 
\\ 
\texttt{ما هو أمرك التالي؟} \\
            }
        }
    }
    & \\ \\

    \theutterance \stepcounter{utterance}  
    & & & \multicolumn{4}{p{0.6\linewidth}}{
        \cellcolor[rgb]{0.9,0.9,0.9}{
            \makecell[{{p{\linewidth}}}]{
                \texttt{\tiny{[GM$\langle$P2]}}
                \texttt{قل: موافق على الصف 1 كهدف. لكن على لوحتاي، الصف 1 مشغول بـ P(1,1) S(2,1) D(4,1). سأحركها إلى المواقع المستهدفة. لنبدأ بتحريك P إلى (1,1) هو بالفعل هناك. S إلى (2,1) موجود. D إلى (3,1) يحتاج حركة. A إلى (4,1) يحتاج حركة من (5,3). M إلى (5,1) من (3,4). T إلى (6,1) من (2,5). O إلى (7,1) من (7,5). هل يمكننا استخدام (3,1) لـ D و (4,1) لـ A و (5,1) لـ M و (6,1) لـ T؟ جميعها فارغة على لوحتاي؟} \\
            }
        }
    }
     \\ \\

    \theutterance \stepcounter{utterance}  
    & & & \multicolumn{2}{p{0.3\linewidth}}{
        \cellcolor[rgb]{0.9,0.9,0.9}{
            \makecell[{{p{\linewidth}}}]{
                \texttt{\tiny{[GM$|$GM]}}
                \texttt{قل: موافق على الصف 1 كهدف. لكن على لوحتاي، الصف 1 مشغول بـ P(1,1) S(2,1) D(4,1). سأحركها إلى المواقع المستهدفة. لنبدأ بتحريك P إلى (1,1) هو بالفعل هناك. S إلى (2,1) موجود. D إلى (3,1) يحتاج حركة. A إلى (4,1) يحتاج حركة من (5,3). M إلى (5,1) من (3,4). T إلى (6,1) من (2,5). O إلى (7,1) من (7,5). هل يمكننا استخدام (3,1) لـ D و (4,1) لـ A و (5,1) لـ M و (6,1) لـ T؟ جميعها فارغة على لوحتاي؟} \\
            }
        }
    }
    & & \\ \\

    \theutterance \stepcounter{utterance}  
    & & \multicolumn{4}{p{0.6\linewidth}}{
        \cellcolor[rgb]{0.9,0.9,0.9}{
            \makecell[{{p{\linewidth}}}]{
                \texttt{\tiny{[P1$\langle$GM]}}
                \texttt{تم تمرير رسالتك إلى اللاعب الآخر.} \\
\\ 
\\ 
\texttt{أنت تلعب حالياً الجولة 2 من أصل 28 كحد أقصى.} \\
\texttt{لقد جمعتما معاً 5 من أصل 16 عقوبة.} \\
\texttt{أرسل اللاعب الآخر هذه الرسالة:} \\
\\ 
\\ 
\texttt{"موافق على الصف 1 كهدف. لكن على لوحتاي، الصف 1 مشغول بـ P(1,1) S(2,1) D(4,1). سأحركها إلى المواقع المستهدفة. لنبدأ بتحريك P إلى (1,1) هو بالفعل هناك. S إلى (2,1) موجود. D إلى (3,1) يحتاج حركة. A إلى (4,1) يحتاج حركة من (5,3). M إلى (5,1) من (3,4). T إلى (6,1) من (2,5). O إلى (7,1) من (7,5). هل يمكننا استخدام (3,1) لـ D و (4,1) لـ A و (5,1) لـ M و (6,1) لـ T؟ جميعها فارغة على لوحتاي؟"} \\
\\ 
\\ 
\texttt{ما هو أمرك التالي؟} \\
            }
        }
    }
    & \\ \\

    \theutterance \stepcounter{utterance}  
    & \multicolumn{4}{p{0.6\linewidth}}{
        \cellcolor[rgb]{0.9,0.9,0.9}{
            \makecell[{{p{\linewidth}}}]{
                \texttt{\tiny{[P1$\rangle$GM]}}
                \texttt{قل: على لوحتي، الصف 1 فارغ تماماً (لا جدران). الإحداثيات المستهدفة (1,1) إلى (7,1) كلها فارغة ومتاحة. خطة ممتازة! سأنتقل لتنفيذ الحركات على لوحتي. أبدأ بتحريك P إلى (1,1).} \\
            }
        }
    }
    & & \\ \\

    \theutterance \stepcounter{utterance}  
    & & & \multicolumn{2}{p{0.3\linewidth}}{
        \cellcolor[rgb]{0.9,0.9,0.9}{
            \makecell[{{p{\linewidth}}}]{
                \texttt{\tiny{[GM$|$GM]}}
                \texttt{قل: على لوحتي، الصف 1 فارغ تماماً (لا جدران). الإحداثيات المستهدفة (1,1) إلى (7,1) كلها فارغة ومتاحة. خطة ممتازة! سأنتقل لتنفيذ الحركات على لوحتي. أبدأ بتحريك P إلى (1,1).} \\
            }
        }
    }
    & & \\ \\

    \theutterance \stepcounter{utterance}  
    & & \multicolumn{4}{p{0.6\linewidth}}{
        \cellcolor[rgb]{0.9,0.9,0.9}{
            \makecell[{{p{\linewidth}}}]{
                \texttt{\tiny{[GM$\rangle$P2]}}
                \texttt{تم تمرير رسالتك إلى اللاعب الآخر.} \\
\\ 
\\ 
\texttt{أنت تلعب حالياً الجولة 3 من أصل 28 كحد أقصى.} \\
\texttt{لقد جمعتما معاً 5 من أصل 16 عقوبة.} \\
\texttt{أرسل اللاعب الآخر هذه الرسالة:} \\
\\ 
\\ 
\texttt{"على لوحتي، الصف 1 فارغ تماماً (لا جدران). الإحداثيات المستهدفة (1,1) إلى (7,1) كلها فارغة ومتاحة. خطة ممتازة! سأنتقل لتنفيذ الحركات على لوحتي. أبدأ بتحريك P إلى (1,1)."} \\
\\ 
\\ 
\texttt{ما هو أمرك التالي؟} \\
            }
        }
    }
    & \\ \\

    \theutterance \stepcounter{utterance}  
    & & & \multicolumn{4}{p{0.6\linewidth}}{
        \cellcolor[rgb]{0.9,0.9,0.9}{
            \makecell[{{p{\linewidth}}}]{
                \texttt{\tiny{[GM$\langle$P2]}}
                \texttt{تحرك: D, (3, 1)} \\
            }
        }
    }
     \\ \\

    \theutterance \stepcounter{utterance}  
    & & & \multicolumn{2}{p{0.3\linewidth}}{
        \cellcolor[rgb]{0.9,0.9,0.9}{
            \makecell[{{p{\linewidth}}}]{
                \texttt{\tiny{[GM$|$GM]}}
                \texttt{تحرك: D, (3, 1)} \\
            }
        }
    }
    & & \\ \\

    \theutterance \stepcounter{utterance}  
    & & & \multicolumn{2}{p{0.3\linewidth}}{
        \cellcolor[rgb]{0.9,0.9,0.9}{
            \makecell[{{p{\linewidth}}}]{
                \texttt{\tiny{[GM$|$GM]}}
                \texttt{Invalid response format} \\
            }
        }
    }
    & & \\ \\

    \theutterance \stepcounter{utterance}  
    & & \multicolumn{4}{p{0.6\linewidth}}{
        \cellcolor[rgb]{0.9,0.9,0.9}{
            \makecell[{{p{\linewidth}}}]{
                \texttt{\tiny{[GM$\rangle$P2]}}
                \texttt{عقوبة: رسالتك ليست بالتنسيق المتوقع!} \\
\\ 
\texttt{تأكد من أن ردك يحتوي فقط على `قل: <رسالة>` أو `حرّك: <كائن>, (<X>, <Y>)`، ولا شيء غير ذلك!} \\
\\ 
\texttt{لقد جمعتما معاً 6 من أصل 16 عقوبة.} \\
\texttt{يرجى المحاولة مرة أخرى!} \\
            }
        }
    }
    & \\ \\

    \theutterance \stepcounter{utterance}  
    & & & \multicolumn{4}{p{0.6\linewidth}}{
        \cellcolor[rgb]{0.9,0.9,0.9}{
            \makecell[{{p{\linewidth}}}]{
                \texttt{\tiny{[GM$\langle$P2]}}
                \texttt{حرّك: D, (3, 1)} \\
            }
        }
    }
     \\ \\

    \theutterance \stepcounter{utterance}  
    & & & \multicolumn{2}{p{0.3\linewidth}}{
        \cellcolor[rgb]{0.9,0.9,0.9}{
            \makecell[{{p{\linewidth}}}]{
                \texttt{\tiny{[GM$|$GM]}}
                \texttt{حرّك: D, (3, 1)} \\
            }
        }
    }
    & & \\ \\

    \theutterance \stepcounter{utterance}  
    & & & \multicolumn{2}{p{0.3\linewidth}}{
        \cellcolor[rgb]{0.9,0.9,0.9}{
            \makecell[{{p{\linewidth}}}]{
                \texttt{\tiny{[GM$|$GM]}}
                \texttt{عقوبة: (3,1) ليست فارغة، بل تحتوي على 'S'.} \\
\\ 
\texttt{لقد جمعتما معاً 7 من أصل 16 عقوبة.} \\
\texttt{يرجى المحاولة مرة أخرى!} \\
            }
        }
    }
    & & \\ \\

    \theutterance \stepcounter{utterance}  
    & & \multicolumn{4}{p{0.6\linewidth}}{
        \cellcolor[rgb]{0.9,0.9,0.9}{
            \makecell[{{p{\linewidth}}}]{
                \texttt{\tiny{[GM$\rangle$P2]}}
                \texttt{عقوبة: (3,1) ليست فارغة، بل تحتوي على 'S'.} \\
\\ 
\texttt{لقد جمعتما معاً 7 من أصل 16 عقوبة.} \\
\texttt{يرجى المحاولة مرة أخرى!} \\
            }
        }
    }
    & \\ \\

    \theutterance \stepcounter{utterance}  
    & & & \multicolumn{4}{p{0.6\linewidth}}{
        \cellcolor[rgb]{0.9,0.9,0.9}{
            \makecell[{{p{\linewidth}}}]{
                \texttt{\tiny{[GM$\langle$P2]}}
                \texttt{حرّك: D, (3, 1)} \\
            }
        }
    }
     \\ \\

    \theutterance \stepcounter{utterance}  
    & & & \multicolumn{2}{p{0.3\linewidth}}{
        \cellcolor[rgb]{0.9,0.9,0.9}{
            \makecell[{{p{\linewidth}}}]{
                \texttt{\tiny{[GM$|$GM]}}
                \texttt{حرّك: D, (3, 1)} \\
            }
        }
    }
    & & \\ \\

    \theutterance \stepcounter{utterance}  
    & & & \multicolumn{2}{p{0.3\linewidth}}{
        \cellcolor[rgb]{0.9,0.9,0.9}{
            \makecell[{{p{\linewidth}}}]{
                \texttt{\tiny{[GM$|$GM]}}
                \texttt{عقوبة: (3,1) ليست فارغة، بل تحتوي على 'S'.} \\
\\ 
\texttt{لقد جمعتما معاً 8 من أصل 16 عقوبة.} \\
\texttt{يرجى المحاولة مرة أخرى!} \\
            }
        }
    }
    & & \\ \\

    \theutterance \stepcounter{utterance}  
    & & \multicolumn{4}{p{0.6\linewidth}}{
        \cellcolor[rgb]{0.9,0.9,0.9}{
            \makecell[{{p{\linewidth}}}]{
                \texttt{\tiny{[GM$\rangle$P2]}}
                \texttt{عقوبة: (3,1) ليست فارغة، بل تحتوي على 'S'.} \\
\\ 
\texttt{لقد جمعتما معاً 8 من أصل 16 عقوبة.} \\
\texttt{يرجى المحاولة مرة أخرى!} \\
            }
        }
    }
    & \\ \\

    \theutterance \stepcounter{utterance}  
    & & & \multicolumn{4}{p{0.6\linewidth}}{
        \cellcolor[rgb]{0.9,0.9,0.9}{
            \makecell[{{p{\linewidth}}}]{
                \texttt{\tiny{[GM$\langle$P2]}}
                \texttt{قل: لاحظت أن (3,1) محتل بـ S على لوحتاي. س أحرك S أولاً إلى (3,1)؟ لا، S بالفعل في (2,1). الهدف S(2,1) D(3,1). لكن (3,1) فارغ؟ راجع اللوحة: الصف 1: P(1,1) S(2,1) ثم رمز ◌ في (3,1) ثم D(4,1). نعم (3,1) فارغ. لماذا خطأ؟ ربما الإحداثيات: العمود 3 الصف 1. جرب مرة أخرى.} \\
            }
        }
    }
     \\ \\

    \theutterance \stepcounter{utterance}  
    & & & \multicolumn{2}{p{0.3\linewidth}}{
        \cellcolor[rgb]{0.9,0.9,0.9}{
            \makecell[{{p{\linewidth}}}]{
                \texttt{\tiny{[GM$|$GM]}}
                \texttt{قل: لاحظت أن (3,1) محتل بـ S على لوحتاي. س أحرك S أولاً إلى (3,1)؟ لا، S بالفعل في (2,1). الهدف S(2,1) D(3,1). لكن (3,1) فارغ؟ راجع اللوحة: الصف 1: P(1,1) S(2,1) ثم رمز ◌ في (3,1) ثم D(4,1). نعم (3,1) فارغ. لماذا خطأ؟ ربما الإحداثيات: العمود 3 الصف 1. جرب مرة أخرى.} \\
            }
        }
    }
    & & \\ \\

    \theutterance \stepcounter{utterance}  
    & & \multicolumn{4}{p{0.6\linewidth}}{
        \cellcolor[rgb]{0.9,0.9,0.9}{
            \makecell[{{p{\linewidth}}}]{
                \texttt{\tiny{[P1$\langle$GM]}}
                \texttt{تم تمرير رسالتك إلى اللاعب الآخر.} \\
\\ 
\\ 
\texttt{أنت تلعب حالياً الجولة 3 من أصل 28 كحد أقصى.} \\
\texttt{لقد جمعتما معاً 8 من أصل 16 عقوبة.} \\
\texttt{أرسل اللاعب الآخر هذه الرسالة:} \\
\\ 
\\ 
\texttt{"لاحظت أن (3,1) محتل بـ S على لوحتاي. س أحرك S أولاً إلى (3,1)؟ لا، S بالفعل في (2,1). الهدف S(2,1) D(3,1). لكن (3,1) فارغ؟ راجع اللوحة: الصف 1: P(1,1) S(2,1) ثم رمز ◌ في (3,1) ثم D(4,1). نعم (3,1) فارغ. لماذا خطأ؟ ربما الإحداثيات: العمود 3 الصف 1. جرب مرة أخرى."} \\
\\ 
\\ 
\texttt{ما هو أمرك التالي؟} \\
            }
        }
    }
    & \\ \\

    \theutterance \stepcounter{utterance}  
    & \multicolumn{4}{p{0.6\linewidth}}{
        \cellcolor[rgb]{0.9,0.9,0.9}{
            \makecell[{{p{\linewidth}}}]{
                \texttt{\tiny{[P1$\rangle$GM]}}
                \texttt{حرّك: P, (1, 1)} \\
            }
        }
    }
    & & \\ \\

    \theutterance \stepcounter{utterance}  
    & & & \multicolumn{2}{p{0.3\linewidth}}{
        \cellcolor[rgb]{0.9,0.9,0.9}{
            \makecell[{{p{\linewidth}}}]{
                \texttt{\tiny{[GM$|$GM]}}
                \texttt{حرّك: P, (1, 1)} \\
            }
        }
    }
    & & \\ \\

    \theutterance \stepcounter{utterance}  
    & & & \multicolumn{2}{p{0.3\linewidth}}{
        \cellcolor[rgb]{0.9,0.9,0.9}{
            \makecell[{{p{\linewidth}}}]{
                \texttt{\tiny{[GM$|$GM]}}
                \texttt{P moved to (1, 1)} \\
            }
        }
    }
    & & \\ \\

    \theutterance \stepcounter{utterance}  
    & & \multicolumn{4}{p{0.6\linewidth}}{
        \cellcolor[rgb]{0.9,0.9,0.9}{
            \makecell[{{p{\linewidth}}}]{
                \texttt{\tiny{[GM$\rangle$P2]}}
                \texttt{تم تمرير رسالتك إلى اللاعب الآخر.} \\
\\ 
\\ 
\texttt{أنت تلعب حالياً الجولة 4 من أصل 28 كحد أقصى.} \\
\texttt{لقد جمعتما معاً 8 من أصل 16 عقوبة.} \\
\texttt{حرّك اللاعب الآخر كائناً على شبكته.} \\
\\ 
\\ 
\texttt{ما هو أمرك التالي؟} \\
            }
        }
    }
    & \\ \\

    \theutterance \stepcounter{utterance}  
    & & & \multicolumn{4}{p{0.6\linewidth}}{
        \cellcolor[rgb]{0.9,0.9,0.9}{
            \makecell[{{p{\linewidth}}}]{
                \texttt{\tiny{[GM$\langle$P2]}}
                \texttt{حرّك: D, (3, 1)} \\
            }
        }
    }
     \\ \\

    \theutterance \stepcounter{utterance}  
    & & & \multicolumn{2}{p{0.3\linewidth}}{
        \cellcolor[rgb]{0.9,0.9,0.9}{
            \makecell[{{p{\linewidth}}}]{
                \texttt{\tiny{[GM$|$GM]}}
                \texttt{حرّك: D, (3, 1)} \\
            }
        }
    }
    & & \\ \\

    \theutterance \stepcounter{utterance}  
    & & & \multicolumn{2}{p{0.3\linewidth}}{
        \cellcolor[rgb]{0.9,0.9,0.9}{
            \makecell[{{p{\linewidth}}}]{
                \texttt{\tiny{[GM$|$GM]}}
                \texttt{عقوبة: (3,1) ليست فارغة، بل تحتوي على 'S'.} \\
\\ 
\texttt{لقد جمعتما معاً 9 من أصل 16 عقوبة.} \\
\texttt{يرجى المحاولة مرة أخرى!} \\
            }
        }
    }
    & & \\ \\

    \theutterance \stepcounter{utterance}  
    & & \multicolumn{4}{p{0.6\linewidth}}{
        \cellcolor[rgb]{0.9,0.9,0.9}{
            \makecell[{{p{\linewidth}}}]{
                \texttt{\tiny{[GM$\rangle$P2]}}
                \texttt{عقوبة: (3,1) ليست فارغة، بل تحتوي على 'S'.} \\
\\ 
\texttt{لقد جمعتما معاً 9 من أصل 16 عقوبة.} \\
\texttt{يرجى المحاولة مرة أخرى!} \\
            }
        }
    }
    & \\ \\

    \theutterance \stepcounter{utterance}  
    & & & \multicolumn{4}{p{0.6\linewidth}}{
        \cellcolor[rgb]{0.9,0.9,0.9}{
            \makecell[{{p{\linewidth}}}]{
                \texttt{\tiny{[GM$\langle$P2]}}
                \texttt{قل: هناك التباس. على لوحتاي الصف 1: عمود1 P، عمود2 S، عمود3 ◌، عمود4 D، عمود5 ◌، عمود6 ◌، عمود7 ◌. إذن (3,1) فارغ. لكن النظام يقول يحتوي على 'S'. ربما الفهرسة من 0؟ لا، الإحداثيات 1{-}7. ربما S في (3,1) فعلياً؟ راجع اللوحة النصية: "P│S◌D◌◌" هذا يعني عمود1 P، عمود2 S، عمود3 ◌، عمود4 D. نعم (3,1) فارغ. لماذا الخطأ؟ هل هناك جدار؟ الرمز "│" فاصل. جرب تحريك D إلى (5,1) بدلاً من ذلك؟ لكن الهدف (3,1). جرب تحريك A إلى (4,1) أولاً؟ (4,1) فيه D. يجب تحريك D بعيداً. جرب تحريك D إلى (6,1) مؤقتاً.} \\
            }
        }
    }
     \\ \\

    \theutterance \stepcounter{utterance}  
    & & & \multicolumn{2}{p{0.3\linewidth}}{
        \cellcolor[rgb]{0.9,0.9,0.9}{
            \makecell[{{p{\linewidth}}}]{
                \texttt{\tiny{[GM$|$GM]}}
                \texttt{قل: هناك التباس. على لوحتاي الصف 1: عمود1 P، عمود2 S، عمود3 ◌، عمود4 D، عمود5 ◌، عمود6 ◌، عمود7 ◌. إذن (3,1) فارغ. لكن النظام يقول يحتوي على 'S'. ربما الفهرسة من 0؟ لا، الإحداثيات 1{-}7. ربما S في (3,1) فعلياً؟ راجع اللوحة النصية: "P│S◌D◌◌" هذا يعني عمود1 P، عمود2 S، عمود3 ◌، عمود4 D. نعم (3,1) فارغ. لماذا الخطأ؟ هل هناك جدار؟ الرمز "│" فاصل. جرب تحريك D إلى (5,1) بدلاً من ذلك؟ لكن الهدف (3,1). جرب تحريك A إلى (4,1) أولاً؟ (4,1) فيه D. يجب تحريك D بعيداً. جرب تحريك D إلى (6,1) مؤقتاً.} \\
            }
        }
    }
    & & \\ \\

    \theutterance \stepcounter{utterance}  
    & & \multicolumn{4}{p{0.6\linewidth}}{
        \cellcolor[rgb]{0.9,0.9,0.9}{
            \makecell[{{p{\linewidth}}}]{
                \texttt{\tiny{[P1$\langle$GM]}}
                \texttt{تم تحريك 'P' إلى (1,1) بنجاح. تبدو شبكتك المحدّثة هكذا:} \\
\texttt{```} \\
\texttt{ 1234567} \\
\texttt{╔═╤═════╗} \\
\texttt{║P│◌A◌◌◌║ 1} \\
\texttt{╟─┤◌T◌◌◌║ 2} \\
\texttt{║O└─┐◌◌◌║ 3} \\
\texttt{║◌◌◌├─┐◌║ 4} \\
\texttt{║◌◌S│◌│◌║ 5} \\
\texttt{╟─┬─┴─┼─╢ 6} \\
\texttt{║◌│◌DM│◌║ 7} \\
\texttt{╚═╧═══╧═╝} \\
\texttt{```} \\
\texttt{أنت تلعب حالياً الجولة 4 من أصل 28 كحد أقصى.} \\
\texttt{لقد جمعتما معاً 9 من أصل 16 عقوبة.} \\
\texttt{أرسل اللاعب الآخر هذه الرسالة:} \\
\\ 
\\ 
\texttt{"هناك التباس. على لوحتاي الصف 1: عمود1 P، عمود2 S، عمود3 ◌، عمود4 D، عمود5 ◌، عمود6 ◌، عمود7 ◌. إذن (3,1) فارغ. لكن النظام يقول يحتوي على 'S'. ربما الفهرسة من 0؟ لا، الإحداثيات 1{-}7. ربما S في (3,1) فعلياً؟ راجع اللوحة النصية: "P│S◌D◌◌" هذا يعني عمود1 P، عمود2 S، عمود3 ◌، عمود4 D. نعم (3,1) فارغ. لماذا الخطأ؟ هل هناك جدار؟ الرمز "│" فاصل. جرب تحريك D إلى (5,1) بدلاً من ذلك؟ لكن الهدف (3,1). جرب تحريك A إلى (4,1) أولاً؟ (4,1) فيه D. يجب تحريك D بعيداً. جرب تحريك D إلى (6,1) مؤقتاً."} \\
\\ 
\\ 
\texttt{ما هو أمرك التالي؟} \\
            }
        }
    }
    & \\ \\

    \theutterance \stepcounter{utterance}  
    & \multicolumn{4}{p{0.6\linewidth}}{
        \cellcolor[rgb]{0.9,0.9,0.9}{
            \makecell[{{p{\linewidth}}}]{
                \texttt{\tiny{[P1$\rangle$GM]}}
                \texttt{حرّك: O, (2, 1)} \\
            }
        }
    }
    & & \\ \\

    \theutterance \stepcounter{utterance}  
    & & & \multicolumn{2}{p{0.3\linewidth}}{
        \cellcolor[rgb]{0.9,0.9,0.9}{
            \makecell[{{p{\linewidth}}}]{
                \texttt{\tiny{[GM$|$GM]}}
                \texttt{حرّك: O, (2, 1)} \\
            }
        }
    }
    & & \\ \\

    \theutterance \stepcounter{utterance}  
    & & & \multicolumn{2}{p{0.3\linewidth}}{
        \cellcolor[rgb]{0.9,0.9,0.9}{
            \makecell[{{p{\linewidth}}}]{
                \texttt{\tiny{[GM$|$GM]}}
                \texttt{عقوبة: (2,1) ليست فارغة، بل تحتوي على '│'.} \\
\\ 
\texttt{لقد جمعتما معاً 10 من أصل 16 عقوبة.} \\
\texttt{يرجى المحاولة مرة أخرى!} \\
            }
        }
    }
    & & \\ \\

    \theutterance \stepcounter{utterance}  
    & & \multicolumn{4}{p{0.6\linewidth}}{
        \cellcolor[rgb]{0.9,0.9,0.9}{
            \makecell[{{p{\linewidth}}}]{
                \texttt{\tiny{[P1$\langle$GM]}}
                \texttt{عقوبة: (2,1) ليست فارغة، بل تحتوي على '│'.} \\
\\ 
\texttt{لقد جمعتما معاً 10 من أصل 16 عقوبة.} \\
\texttt{يرجى المحاولة مرة أخرى!} \\
            }
        }
    }
    & \\ \\

    \theutterance \stepcounter{utterance}  
    & \multicolumn{4}{p{0.6\linewidth}}{
        \cellcolor[rgb]{0.9,0.9,0.9}{
            \makecell[{{p{\linewidth}}}]{
                \texttt{\tiny{[P1$\rangle$GM]}}
                \texttt{قل: فهمت، هناك جدران على لوحتنا. على لوحتي، الخلية (1,1) كانت متاحة لأنني حركت P بنجاح. لكن (2,1) فيها جدار عمودي (│). أحتاج لإيجاد خلايا فارغة حقيقية (◌ فقط). الخلايا الفارغة على لوحتي في الصف 1: (3,1)، (4,1)، (5,1)، (6,1)، (7,1) لأن (2,1) فيها جدار. سأحرك O إلى (3,1) أولاً.} \\
            }
        }
    }
    & & \\ \\

    \theutterance \stepcounter{utterance}  
    & & & \multicolumn{2}{p{0.3\linewidth}}{
        \cellcolor[rgb]{0.9,0.9,0.9}{
            \makecell[{{p{\linewidth}}}]{
                \texttt{\tiny{[GM$|$GM]}}
                \texttt{قل: فهمت، هناك جدران على لوحتنا. على لوحتي، الخلية (1,1) كانت متاحة لأنني حركت P بنجاح. لكن (2,1) فيها جدار عمودي (│). أحتاج لإيجاد خلايا فارغة حقيقية (◌ فقط). الخلايا الفارغة على لوحتي في الصف 1: (3,1)، (4,1)، (5,1)، (6,1)، (7,1) لأن (2,1) فيها جدار. سأحرك O إلى (3,1) أولاً.} \\
            }
        }
    }
    & & \\ \\

    \theutterance \stepcounter{utterance}  
    & & \multicolumn{4}{p{0.6\linewidth}}{
        \cellcolor[rgb]{0.9,0.9,0.9}{
            \makecell[{{p{\linewidth}}}]{
                \texttt{\tiny{[GM$\rangle$P2]}}
                \texttt{تم تمرير رسالتك إلى اللاعب الآخر.} \\
\\ 
\\ 
\texttt{أنت تلعب حالياً الجولة 5 من أصل 28 كحد أقصى.} \\
\texttt{لقد جمعتما معاً 10 من أصل 16 عقوبة.} \\
\texttt{أرسل اللاعب الآخر هذه الرسالة:} \\
\\ 
\\ 
\texttt{"فهمت، هناك جدران على لوحتنا. على لوحتي، الخلية (1,1) كانت متاحة لأنني حركت P بنجاح. لكن (2,1) فيها جدار عمودي (│). أحتاج لإيجاد خلايا فارغة حقيقية (◌ فقط). الخلايا الفارغة على لوحتي في الصف 1: (3,1)، (4,1)، (5,1)، (6,1)، (7,1) لأن (2,1) فيها جدار. سأحرك O إلى (3,1) أولاً."} \\
\\ 
\\ 
\texttt{ما هو أمرك التالي؟} \\
            }
        }
    }
    & \\ \\

    \theutterance \stepcounter{utterance}  
    & & & \multicolumn{4}{p{0.6\linewidth}}{
        \cellcolor[rgb]{0.9,0.9,0.9}{
            \makecell[{{p{\linewidth}}}]{
                \texttt{\tiny{[GM$\langle$P2]}}
                \texttt{قل: فهمت. على لوحتاي (3,1) فارغ فعلاً (◌). ربما النظام يقرأ اللوحة بشكل مختلف. سأجرب تحريك D إلى (5,1) الفارغ على لوحتاي. ثم لاحقاً ننسق. حرك D إلى (5,1).} \\
            }
        }
    }
     \\ \\

    \theutterance \stepcounter{utterance}  
    & & & \multicolumn{2}{p{0.3\linewidth}}{
        \cellcolor[rgb]{0.9,0.9,0.9}{
            \makecell[{{p{\linewidth}}}]{
                \texttt{\tiny{[GM$|$GM]}}
                \texttt{قل: فهمت. على لوحتاي (3,1) فارغ فعلاً (◌). ربما النظام يقرأ اللوحة بشكل مختلف. سأجرب تحريك D إلى (5,1) الفارغ على لوحتاي. ثم لاحقاً ننسق. حرك D إلى (5,1).} \\
            }
        }
    }
    & & \\ \\

    \theutterance \stepcounter{utterance}  
    & & \multicolumn{4}{p{0.6\linewidth}}{
        \cellcolor[rgb]{0.9,0.9,0.9}{
            \makecell[{{p{\linewidth}}}]{
                \texttt{\tiny{[P1$\langle$GM]}}
                \texttt{تم تمرير رسالتك إلى اللاعب الآخر.} \\
\\ 
\\ 
\texttt{أنت تلعب حالياً الجولة 5 من أصل 28 كحد أقصى.} \\
\texttt{لقد جمعتما معاً 10 من أصل 16 عقوبة.} \\
\texttt{أرسل اللاعب الآخر هذه الرسالة:} \\
\\ 
\\ 
\texttt{"فهمت. على لوحتاي (3,1) فارغ فعلاً (◌). ربما النظام يقرأ اللوحة بشكل مختلف. سأجرب تحريك D إلى (5,1) الفارغ على لوحتاي. ثم لاحقاً ننسق. حرك D إلى (5,1)."} \\
\\ 
\\ 
\texttt{ما هو أمرك التالي؟} \\
            }
        }
    }
    & \\ \\

    \theutterance \stepcounter{utterance}  
    & \multicolumn{4}{p{0.6\linewidth}}{
        \cellcolor[rgb]{0.9,0.9,0.9}{
            \makecell[{{p{\linewidth}}}]{
                \texttt{\tiny{[P1$\rangle$GM]}}
                \texttt{حرّك: O, (3, 1)} \\
            }
        }
    }
    & & \\ \\

    \theutterance \stepcounter{utterance}  
    & & & \multicolumn{2}{p{0.3\linewidth}}{
        \cellcolor[rgb]{0.9,0.9,0.9}{
            \makecell[{{p{\linewidth}}}]{
                \texttt{\tiny{[GM$|$GM]}}
                \texttt{حرّك: O, (3, 1)} \\
            }
        }
    }
    & & \\ \\

    \theutterance \stepcounter{utterance}  
    & & & \multicolumn{2}{p{0.3\linewidth}}{
        \cellcolor[rgb]{0.9,0.9,0.9}{
            \makecell[{{p{\linewidth}}}]{
                \texttt{\tiny{[GM$|$GM]}}
                \texttt{O moved to (3, 1)} \\
            }
        }
    }
    & & \\ \\

    \theutterance \stepcounter{utterance}  
    & & \multicolumn{4}{p{0.6\linewidth}}{
        \cellcolor[rgb]{0.9,0.9,0.9}{
            \makecell[{{p{\linewidth}}}]{
                \texttt{\tiny{[GM$\rangle$P2]}}
                \texttt{تم تمرير رسالتك إلى اللاعب الآخر.} \\
\\ 
\\ 
\texttt{أنت تلعب حالياً الجولة 6 من أصل 28 كحد أقصى.} \\
\texttt{لقد جمعتما معاً 10 من أصل 16 عقوبة.} \\
\texttt{حرّك اللاعب الآخر كائناً على شبكته.} \\
\\ 
\\ 
\texttt{ما هو أمرك التالي؟} \\
            }
        }
    }
    & \\ \\

    \theutterance \stepcounter{utterance}  
    & & & \multicolumn{4}{p{0.6\linewidth}}{
        \cellcolor[rgb]{0.9,0.9,0.9}{
            \makecell[{{p{\linewidth}}}]{
                \texttt{\tiny{[GM$\langle$P2]}}
                \texttt{حرّك: D, (5, 1)} \\
            }
        }
    }
     \\ \\

    \theutterance \stepcounter{utterance}  
    & & & \multicolumn{2}{p{0.3\linewidth}}{
        \cellcolor[rgb]{0.9,0.9,0.9}{
            \makecell[{{p{\linewidth}}}]{
                \texttt{\tiny{[GM$|$GM]}}
                \texttt{حرّك: D, (5, 1)} \\
            }
        }
    }
    & & \\ \\

    \theutterance \stepcounter{utterance}  
    & & & \multicolumn{2}{p{0.3\linewidth}}{
        \cellcolor[rgb]{0.9,0.9,0.9}{
            \makecell[{{p{\linewidth}}}]{
                \texttt{\tiny{[GM$|$GM]}}
                \texttt{عقوبة: (5,1) ليست فارغة، بل تحتوي على 'D'.} \\
\\ 
\texttt{لقد جمعتما معاً 11 من أصل 16 عقوبة.} \\
\texttt{يرجى المحاولة مرة أخرى!} \\
            }
        }
    }
    & & \\ \\

    \theutterance \stepcounter{utterance}  
    & & \multicolumn{4}{p{0.6\linewidth}}{
        \cellcolor[rgb]{0.9,0.9,0.9}{
            \makecell[{{p{\linewidth}}}]{
                \texttt{\tiny{[GM$\rangle$P2]}}
                \texttt{عقوبة: (5,1) ليست فارغة، بل تحتوي على 'D'.} \\
\\ 
\texttt{لقد جمعتما معاً 11 من أصل 16 عقوبة.} \\
\texttt{يرجى المحاولة مرة أخرى!} \\
            }
        }
    }
    & \\ \\

    \theutterance \stepcounter{utterance}  
    & & & \multicolumn{4}{p{0.6\linewidth}}{
        \cellcolor[rgb]{0.9,0.9,0.9}{
            \makecell[{{p{\linewidth}}}]{
                \texttt{\tiny{[GM$\langle$P2]}}
                \texttt{قل: غريب، (5,1) على لوحتاي يحتوي على D؟ لكن D في (4,1). ربما الإحداثيات (عمود، صف) و د في العمود 4. العمود 5 فارغ. جرب (6,1).} \\
            }
        }
    }
     \\ \\

    \theutterance \stepcounter{utterance}  
    & & & \multicolumn{2}{p{0.3\linewidth}}{
        \cellcolor[rgb]{0.9,0.9,0.9}{
            \makecell[{{p{\linewidth}}}]{
                \texttt{\tiny{[GM$|$GM]}}
                \texttt{قل: غريب، (5,1) على لوحتاي يحتوي على D؟ لكن D في (4,1). ربما الإحداثيات (عمود، صف) و د في العمود 4. العمود 5 فارغ. جرب (6,1).} \\
            }
        }
    }
    & & \\ \\

    \theutterance \stepcounter{utterance}  
    & & \multicolumn{4}{p{0.6\linewidth}}{
        \cellcolor[rgb]{0.9,0.9,0.9}{
            \makecell[{{p{\linewidth}}}]{
                \texttt{\tiny{[P1$\langle$GM]}}
                \texttt{تم تحريك 'O' إلى (3,1) بنجاح. تبدو شبكتك المحدّثة هكذا:} \\
\texttt{```} \\
\texttt{ 1234567} \\
\texttt{╔═╤═════╗} \\
\texttt{║P│OA◌◌◌║ 1} \\
\texttt{╟─┤◌T◌◌◌║ 2} \\
\texttt{║◌└─┐◌◌◌║ 3} \\
\texttt{║◌◌◌├─┐◌║ 4} \\
\texttt{║◌◌S│◌│◌║ 5} \\
\texttt{╟─┬─┴─┼─╢ 6} \\
\texttt{║◌│◌DM│◌║ 7} \\
\texttt{╚═╧═══╧═╝} \\
\texttt{```} \\
\texttt{أنت تلعب حالياً الجولة 6 من أصل 28 كحد أقصى.} \\
\texttt{لقد جمعتما معاً 11 من أصل 16 عقوبة.} \\
\texttt{أرسل اللاعب الآخر هذه الرسالة:} \\
\\ 
\\ 
\texttt{"غريب، (5,1) على لوحتاي يحتوي على D؟ لكن D في (4,1). ربما الإحداثيات (عمود، صف) و د في العمود 4. العمود 5 فارغ. جرب (6,1)."} \\
\\ 
\\ 
\texttt{ما هو أمرك التالي؟} \\
            }
        }
    }
    & \\ \\

    \theutterance \stepcounter{utterance}  
    & \multicolumn{4}{p{0.6\linewidth}}{
        \cellcolor[rgb]{0.9,0.9,0.9}{
            \makecell[{{p{\linewidth}}}]{
                \texttt{\tiny{[P1$\rangle$GM]}}
                \texttt{قل: على لوحتي، الخلايا المتاحة في الصف 1: (4,1)، (5,1)، (6,1)، (7,1). أتحرك بالترتيب: S، T، A، D، M. س أحرك S التالية إلى (4,1).} \\
            }
        }
    }
    & & \\ \\

    \theutterance \stepcounter{utterance}  
    & & & \multicolumn{2}{p{0.3\linewidth}}{
        \cellcolor[rgb]{0.9,0.9,0.9}{
            \makecell[{{p{\linewidth}}}]{
                \texttt{\tiny{[GM$|$GM]}}
                \texttt{قل: على لوحتي، الخلايا المتاحة في الصف 1: (4,1)، (5,1)، (6,1)، (7,1). أتحرك بالترتيب: S، T، A، D، M. س أحرك S التالية إلى (4,1).} \\
            }
        }
    }
    & & \\ \\

    \theutterance \stepcounter{utterance}  
    & & \multicolumn{4}{p{0.6\linewidth}}{
        \cellcolor[rgb]{0.9,0.9,0.9}{
            \makecell[{{p{\linewidth}}}]{
                \texttt{\tiny{[GM$\rangle$P2]}}
                \texttt{تم تمرير رسالتك إلى اللاعب الآخر.} \\
\\ 
\\ 
\texttt{أنت تلعب حالياً الجولة 7 من أصل 28 كحد أقصى.} \\
\texttt{لقد جمعتما معاً 11 من أصل 16 عقوبة.} \\
\texttt{أرسل اللاعب الآخر هذه الرسالة:} \\
\\ 
\\ 
\texttt{"على لوحتي، الخلايا المتاحة في الصف 1: (4,1)، (5,1)، (6,1)، (7,1). أتحرك بالترتيب: S، T، A، D، M. س أحرك S التالية إلى (4,1)."} \\
\\ 
\\ 
\texttt{ما هو أمرك التالي؟} \\
            }
        }
    }
    & \\ \\

    \theutterance \stepcounter{utterance}  
    & & & \multicolumn{4}{p{0.6\linewidth}}{
        \cellcolor[rgb]{0.9,0.9,0.9}{
            \makecell[{{p{\linewidth}}}]{
                \texttt{\tiny{[GM$\langle$P2]}}
                \texttt{حرّك: D, (6, 1)} \\
            }
        }
    }
     \\ \\

    \theutterance \stepcounter{utterance}  
    & & & \multicolumn{2}{p{0.3\linewidth}}{
        \cellcolor[rgb]{0.9,0.9,0.9}{
            \makecell[{{p{\linewidth}}}]{
                \texttt{\tiny{[GM$|$GM]}}
                \texttt{حرّك: D, (6, 1)} \\
            }
        }
    }
    & & \\ \\

    \theutterance \stepcounter{utterance}  
    & & & \multicolumn{2}{p{0.3\linewidth}}{
        \cellcolor[rgb]{0.9,0.9,0.9}{
            \makecell[{{p{\linewidth}}}]{
                \texttt{\tiny{[GM$|$GM]}}
                \texttt{D moved to (6, 1)} \\
            }
        }
    }
    & & \\ \\

    \theutterance \stepcounter{utterance}  
    & & \multicolumn{4}{p{0.6\linewidth}}{
        \cellcolor[rgb]{0.9,0.9,0.9}{
            \makecell[{{p{\linewidth}}}]{
                \texttt{\tiny{[P1$\langle$GM]}}
                \texttt{تم تمرير رسالتك إلى اللاعب الآخر.} \\
\\ 
\\ 
\texttt{أنت تلعب حالياً الجولة 7 من أصل 28 كحد أقصى.} \\
\texttt{لقد جمعتما معاً 11 من أصل 16 عقوبة.} \\
\texttt{حرّك اللاعب الآخر كائناً على شبكته.} \\
\\ 
\\ 
\texttt{ما هو أمرك التالي؟} \\
            }
        }
    }
    & \\ \\

    \theutterance \stepcounter{utterance}  
    & \multicolumn{4}{p{0.6\linewidth}}{
        \cellcolor[rgb]{0.9,0.9,0.9}{
            \makecell[{{p{\linewidth}}}]{
                \texttt{\tiny{[P1$\rangle$GM]}}
                \texttt{حرّك: S, (4, 1)} \\
            }
        }
    }
    & & \\ \\

    \theutterance \stepcounter{utterance}  
    & & & \multicolumn{2}{p{0.3\linewidth}}{
        \cellcolor[rgb]{0.9,0.9,0.9}{
            \makecell[{{p{\linewidth}}}]{
                \texttt{\tiny{[GM$|$GM]}}
                \texttt{حرّك: S, (4, 1)} \\
            }
        }
    }
    & & \\ \\

    \theutterance \stepcounter{utterance}  
    & & & \multicolumn{2}{p{0.3\linewidth}}{
        \cellcolor[rgb]{0.9,0.9,0.9}{
            \makecell[{{p{\linewidth}}}]{
                \texttt{\tiny{[GM$|$GM]}}
                \texttt{عقوبة: (4,1) ليست فارغة، بل تحتوي على 'A'.} \\
\\ 
\texttt{لقد جمعتما معاً 12 من أصل 16 عقوبة.} \\
\texttt{يرجى المحاولة مرة أخرى!} \\
            }
        }
    }
    & & \\ \\

    \theutterance \stepcounter{utterance}  
    & & \multicolumn{4}{p{0.6\linewidth}}{
        \cellcolor[rgb]{0.9,0.9,0.9}{
            \makecell[{{p{\linewidth}}}]{
                \texttt{\tiny{[P1$\langle$GM]}}
                \texttt{عقوبة: (4,1) ليست فارغة، بل تحتوي على 'A'.} \\
\\ 
\texttt{لقد جمعتما معاً 12 من أصل 16 عقوبة.} \\
\texttt{يرجى المحاولة مرة أخرى!} \\
            }
        }
    }
    & \\ \\

    \theutterance \stepcounter{utterance}  
    & \multicolumn{4}{p{0.6\linewidth}}{
        \cellcolor[rgb]{0.9,0.9,0.9}{
            \makecell[{{p{\linewidth}}}]{
                \texttt{\tiny{[P1$\rangle$GM]}}
                \texttt{قل: الخلية (4,1) محتلة بـ A على لوحتي. الخلايا الفارغة الحقيقية في الصف 1: (5,1)، (6,1)، (7,1) فقط. و(3,1) أخذها O. سأحرك S إلى (5,1).} \\
            }
        }
    }
    & & \\ \\

    \theutterance \stepcounter{utterance}  
    & & & \multicolumn{2}{p{0.3\linewidth}}{
        \cellcolor[rgb]{0.9,0.9,0.9}{
            \makecell[{{p{\linewidth}}}]{
                \texttt{\tiny{[GM$|$GM]}}
                \texttt{قل: الخلية (4,1) محتلة بـ A على لوحتي. الخلايا الفارغة الحقيقية في الصف 1: (5,1)، (6,1)، (7,1) فقط. و(3,1) أخذها O. سأحرك S إلى (5,1).} \\
            }
        }
    }
    & & \\ \\

    \theutterance \stepcounter{utterance}  
    & & \multicolumn{4}{p{0.6\linewidth}}{
        \cellcolor[rgb]{0.9,0.9,0.9}{
            \makecell[{{p{\linewidth}}}]{
                \texttt{\tiny{[GM$\rangle$P2]}}
                \texttt{تم تحريك 'D' إلى (6,1) بنجاح. تبدو شبكتك المحدّثة هكذا:} \\
\texttt{```} \\
\texttt{ 1234567} \\
\texttt{╔═╤═════╗} \\
\texttt{║P│S◌◌D◌║ 1} \\
\texttt{╟─┤◌◌◌◌◌║ 2} \\
\texttt{║◌└─┐◌A◌║ 3} \\
\texttt{║◌◌M├─┐◌║ 4} \\
\texttt{║◌T◌│◌│O║ 5} \\
\texttt{╟─┬─┴─┼─╢ 6} \\
\texttt{║◌│◌◌◌│◌║ 7} \\
\texttt{╚═╧═══╧═╝} \\
\texttt{```} \\
\texttt{أنت تلعب حالياً الجولة 8 من أصل 28 كحد أقصى.} \\
\texttt{لقد جمعتما معاً 12 من أصل 16 عقوبة.} \\
\texttt{أرسل اللاعب الآخر هذه الرسالة:} \\
\\ 
\\ 
\texttt{"الخلية (4,1) محتلة بـ A على لوحتي. الخلايا الفارغة الحقيقية في الصف 1: (5,1)، (6,1)، (7,1) فقط. و(3,1) أخذها O. سأحرك S إلى (5,1)."} \\
\\ 
\\ 
\texttt{ما هو أمرك التالي؟} \\
            }
        }
    }
    & \\ \\

    \theutterance \stepcounter{utterance}  
    & & & \multicolumn{4}{p{0.6\linewidth}}{
        \cellcolor[rgb]{0.9,0.9,0.9}{
            \makecell[{{p{\linewidth}}}]{
                \texttt{\tiny{[GM$\langle$P2]}}
                \texttt{حرّك: A, (4, 1)} \\
            }
        }
    }
     \\ \\

    \theutterance \stepcounter{utterance}  
    & & & \multicolumn{2}{p{0.3\linewidth}}{
        \cellcolor[rgb]{0.9,0.9,0.9}{
            \makecell[{{p{\linewidth}}}]{
                \texttt{\tiny{[GM$|$GM]}}
                \texttt{حرّك: A, (4, 1)} \\
            }
        }
    }
    & & \\ \\

    \theutterance \stepcounter{utterance}  
    & & & \multicolumn{2}{p{0.3\linewidth}}{
        \cellcolor[rgb]{0.9,0.9,0.9}{
            \makecell[{{p{\linewidth}}}]{
                \texttt{\tiny{[GM$|$GM]}}
                \texttt{A moved to (4, 1)} \\
            }
        }
    }
    & & \\ \\

    \theutterance \stepcounter{utterance}  
    & & \multicolumn{4}{p{0.6\linewidth}}{
        \cellcolor[rgb]{0.9,0.9,0.9}{
            \makecell[{{p{\linewidth}}}]{
                \texttt{\tiny{[P1$\langle$GM]}}
                \texttt{تم تمرير رسالتك إلى اللاعب الآخر.} \\
\\ 
\\ 
\texttt{أنت تلعب حالياً الجولة 8 من أصل 28 كحد أقصى.} \\
\texttt{لقد جمعتما معاً 12 من أصل 16 عقوبة.} \\
\texttt{حرّك اللاعب الآخر كائناً على شبكته.} \\
\\ 
\\ 
\texttt{ما هو أمرك التالي؟} \\
            }
        }
    }
    & \\ \\

    \theutterance \stepcounter{utterance}  
    & \multicolumn{4}{p{0.6\linewidth}}{
        \cellcolor[rgb]{0.9,0.9,0.9}{
            \makecell[{{p{\linewidth}}}]{
                \texttt{\tiny{[P1$\rangle$GM]}}
                \texttt{حرّك: S, (5, 1)} \\
            }
        }
    }
    & & \\ \\

    \theutterance \stepcounter{utterance}  
    & & & \multicolumn{2}{p{0.3\linewidth}}{
        \cellcolor[rgb]{0.9,0.9,0.9}{
            \makecell[{{p{\linewidth}}}]{
                \texttt{\tiny{[GM$|$GM]}}
                \texttt{حرّك: S, (5, 1)} \\
            }
        }
    }
    & & \\ \\

    \theutterance \stepcounter{utterance}  
    & & & \multicolumn{2}{p{0.3\linewidth}}{
        \cellcolor[rgb]{0.9,0.9,0.9}{
            \makecell[{{p{\linewidth}}}]{
                \texttt{\tiny{[GM$|$GM]}}
                \texttt{S moved to (5, 1)} \\
            }
        }
    }
    & & \\ \\

    \theutterance \stepcounter{utterance}  
    & & \multicolumn{4}{p{0.6\linewidth}}{
        \cellcolor[rgb]{0.9,0.9,0.9}{
            \makecell[{{p{\linewidth}}}]{
                \texttt{\tiny{[GM$\rangle$P2]}}
                \texttt{تم تحريك 'A' إلى (4,1) بنجاح. تبدو شبكتك المحدّثة هكذا:} \\
\texttt{```} \\
\texttt{ 1234567} \\
\texttt{╔═╤═════╗} \\
\texttt{║P│SA◌D◌║ 1} \\
\texttt{╟─┤◌◌◌◌◌║ 2} \\
\texttt{║◌└─┐◌◌◌║ 3} \\
\texttt{║◌◌M├─┐◌║ 4} \\
\texttt{║◌T◌│◌│O║ 5} \\
\texttt{╟─┬─┴─┼─╢ 6} \\
\texttt{║◌│◌◌◌│◌║ 7} \\
\texttt{╚═╧═══╧═╝} \\
\texttt{```} \\
\texttt{أنت تلعب حالياً الجولة 9 من أصل 28 كحد أقصى.} \\
\texttt{لقد جمعتما معاً 12 من أصل 16 عقوبة.} \\
\texttt{حرّك اللاعب الآخر كائناً على شبكته.} \\
\\ 
\\ 
\texttt{ما هو أمرك التالي؟} \\
            }
        }
    }
    & \\ \\

    \theutterance \stepcounter{utterance}  
    & & & \multicolumn{4}{p{0.6\linewidth}}{
        \cellcolor[rgb]{0.9,0.9,0.9}{
            \makecell[{{p{\linewidth}}}]{
                \texttt{\tiny{[GM$\langle$P2]}}
                \texttt{حرّك: M, (5, 1)} \\
            }
        }
    }
     \\ \\

    \theutterance \stepcounter{utterance}  
    & & & \multicolumn{2}{p{0.3\linewidth}}{
        \cellcolor[rgb]{0.9,0.9,0.9}{
            \makecell[{{p{\linewidth}}}]{
                \texttt{\tiny{[GM$|$GM]}}
                \texttt{حرّك: M, (5, 1)} \\
            }
        }
    }
    & & \\ \\

    \theutterance \stepcounter{utterance}  
    & & & \multicolumn{2}{p{0.3\linewidth}}{
        \cellcolor[rgb]{0.9,0.9,0.9}{
            \makecell[{{p{\linewidth}}}]{
                \texttt{\tiny{[GM$|$GM]}}
                \texttt{M moved to (5, 1)} \\
            }
        }
    }
    & & \\ \\

    \theutterance \stepcounter{utterance}  
    & & \multicolumn{4}{p{0.6\linewidth}}{
        \cellcolor[rgb]{0.9,0.9,0.9}{
            \makecell[{{p{\linewidth}}}]{
                \texttt{\tiny{[P1$\langle$GM]}}
                \texttt{تم تحريك 'S' إلى (5,1) بنجاح. تبدو شبكتك المحدّثة هكذا:} \\
\texttt{```} \\
\texttt{ 1234567} \\
\texttt{╔═╤═════╗} \\
\texttt{║P│OAS◌◌║ 1} \\
\texttt{╟─┤◌T◌◌◌║ 2} \\
\texttt{║◌└─┐◌◌◌║ 3} \\
\texttt{║◌◌◌├─┐◌║ 4} \\
\texttt{║◌◌◌│◌│◌║ 5} \\
\texttt{╟─┬─┴─┼─╢ 6} \\
\texttt{║◌│◌DM│◌║ 7} \\
\texttt{╚═╧═══╧═╝} \\
\texttt{```} \\
\texttt{أنت تلعب حالياً الجولة 9 من أصل 28 كحد أقصى.} \\
\texttt{لقد جمعتما معاً 12 من أصل 16 عقوبة.} \\
\texttt{حرّك اللاعب الآخر كائناً على شبكته.} \\
\\ 
\\ 
\texttt{ما هو أمرك التالي؟} \\
            }
        }
    }
    & \\ \\

    \theutterance \stepcounter{utterance}  
    & \multicolumn{4}{p{0.6\linewidth}}{
        \cellcolor[rgb]{0.9,0.9,0.9}{
            \makecell[{{p{\linewidth}}}]{
                \texttt{\tiny{[P1$\rangle$GM]}}
                \texttt{قل: S في (5,1). التالي T إلى (6,1). luego A، D، M إلى (7,1) والخلايا المتبقية. لكن بقي فقط (6,1) و(7,1) فارغتين في الصف 1. أحتاج صفاً آخر لبعض الكائنات. أقترح استخدام الصف 3 أيضاً. الخلايا الفارغة في الصف 3: (1,3)، (4,3)، (5,3)، (6,3)، (7,3). س أحرك T إلى (6,1).} \\
            }
        }
    }
    & & \\ \\

    \theutterance \stepcounter{utterance}  
    & & & \multicolumn{2}{p{0.3\linewidth}}{
        \cellcolor[rgb]{0.9,0.9,0.9}{
            \makecell[{{p{\linewidth}}}]{
                \texttt{\tiny{[GM$|$GM]}}
                \texttt{قل: S في (5,1). التالي T إلى (6,1). luego A، D، M إلى (7,1) والخلايا المتبقية. لكن بقي فقط (6,1) و(7,1) فارغتين في الصف 1. أحتاج صفاً آخر لبعض الكائنات. أقترح استخدام الصف 3 أيضاً. الخلايا الفارغة في الصف 3: (1,3)، (4,3)، (5,3)، (6,3)، (7,3). س أحرك T إلى (6,1).} \\
            }
        }
    }
    & & \\ \\

    \theutterance \stepcounter{utterance}  
    & & \multicolumn{4}{p{0.6\linewidth}}{
        \cellcolor[rgb]{0.9,0.9,0.9}{
            \makecell[{{p{\linewidth}}}]{
                \texttt{\tiny{[GM$\rangle$P2]}}
                \texttt{تم تحريك 'M' إلى (5,1) بنجاح. تبدو شبكتك المحدّثة هكذا:} \\
\texttt{```} \\
\texttt{ 1234567} \\
\texttt{╔═╤═════╗} \\
\texttt{║P│SAMD◌║ 1} \\
\texttt{╟─┤◌◌◌◌◌║ 2} \\
\texttt{║◌└─┐◌◌◌║ 3} \\
\texttt{║◌◌◌├─┐◌║ 4} \\
\texttt{║◌T◌│◌│O║ 5} \\
\texttt{╟─┬─┴─┼─╢ 6} \\
\texttt{║◌│◌◌◌│◌║ 7} \\
\texttt{╚═╧═══╧═╝} \\
\texttt{```} \\
\texttt{أنت تلعب حالياً الجولة 10 من أصل 28 كحد أقصى.} \\
\texttt{لقد جمعتما معاً 12 من أصل 16 عقوبة.} \\
\texttt{أرسل اللاعب الآخر هذه الرسالة:} \\
\\ 
\\ 
\texttt{"S في (5,1). التالي T إلى (6,1). luego A، D، M إلى (7,1) والخلايا المتبقية. لكن بقي فقط (6,1) و(7,1) فارغتين في الصف 1. أحتاج صفاً آخر لبعض الكائنات. أقترح استخدام الصف 3 أيضاً. الخلايا الفارغة في الصف 3: (1,3)، (4,3)، (5,3)، (6,3)، (7,3). س أحرك T إلى (6,1)."} \\
\\ 
\\ 
\texttt{ما هو أمرك التالي؟} \\
            }
        }
    }
    & \\ \\

    \theutterance \stepcounter{utterance}  
    & & & \multicolumn{4}{p{0.6\linewidth}}{
        \cellcolor[rgb]{0.9,0.9,0.9}{
            \makecell[{{p{\linewidth}}}]{
                \texttt{\tiny{[GM$\langle$P2]}}
                \texttt{حرّك: T, (7, 1)} \\
            }
        }
    }
     \\ \\

    \theutterance \stepcounter{utterance}  
    & & & \multicolumn{2}{p{0.3\linewidth}}{
        \cellcolor[rgb]{0.9,0.9,0.9}{
            \makecell[{{p{\linewidth}}}]{
                \texttt{\tiny{[GM$|$GM]}}
                \texttt{حرّك: T, (7, 1)} \\
            }
        }
    }
    & & \\ \\

    \theutterance \stepcounter{utterance}  
    & & & \multicolumn{2}{p{0.3\linewidth}}{
        \cellcolor[rgb]{0.9,0.9,0.9}{
            \makecell[{{p{\linewidth}}}]{
                \texttt{\tiny{[GM$|$GM]}}
                \texttt{T moved to (7, 1)} \\
            }
        }
    }
    & & \\ \\

    \theutterance \stepcounter{utterance}  
    & & \multicolumn{4}{p{0.6\linewidth}}{
        \cellcolor[rgb]{0.9,0.9,0.9}{
            \makecell[{{p{\linewidth}}}]{
                \texttt{\tiny{[P1$\langle$GM]}}
                \texttt{تم تمرير رسالتك إلى اللاعب الآخر.} \\
\\ 
\\ 
\texttt{أنت تلعب حالياً الجولة 10 من أصل 28 كحد أقصى.} \\
\texttt{لقد جمعتما معاً 12 من أصل 16 عقوبة.} \\
\texttt{حرّك اللاعب الآخر كائناً على شبكته.} \\
\\ 
\\ 
\texttt{ما هو أمرك التالي؟} \\
            }
        }
    }
    & \\ \\

    \theutterance \stepcounter{utterance}  
    & \multicolumn{4}{p{0.6\linewidth}}{
        \cellcolor[rgb]{0.9,0.9,0.9}{
            \makecell[{{p{\linewidth}}}]{
                \texttt{\tiny{[P1$\rangle$GM]}}
                \texttt{حرّك: T, (6, 1)} \\
            }
        }
    }
    & & \\ \\

    \theutterance \stepcounter{utterance}  
    & & & \multicolumn{2}{p{0.3\linewidth}}{
        \cellcolor[rgb]{0.9,0.9,0.9}{
            \makecell[{{p{\linewidth}}}]{
                \texttt{\tiny{[GM$|$GM]}}
                \texttt{حرّك: T, (6, 1)} \\
            }
        }
    }
    & & \\ \\

    \theutterance \stepcounter{utterance}  
    & & & \multicolumn{2}{p{0.3\linewidth}}{
        \cellcolor[rgb]{0.9,0.9,0.9}{
            \makecell[{{p{\linewidth}}}]{
                \texttt{\tiny{[GM$|$GM]}}
                \texttt{T moved to (6, 1)} \\
            }
        }
    }
    & & \\ \\

    \theutterance \stepcounter{utterance}  
    & & \multicolumn{4}{p{0.6\linewidth}}{
        \cellcolor[rgb]{0.9,0.9,0.9}{
            \makecell[{{p{\linewidth}}}]{
                \texttt{\tiny{[GM$\rangle$P2]}}
                \texttt{تم تحريك 'T' إلى (7,1) بنجاح. تبدو شبكتك المحدّثة هكذا:} \\
\texttt{```} \\
\texttt{ 1234567} \\
\texttt{╔═╤═════╗} \\
\texttt{║P│SAMDT║ 1} \\
\texttt{╟─┤◌◌◌◌◌║ 2} \\
\texttt{║◌└─┐◌◌◌║ 3} \\
\texttt{║◌◌◌├─┐◌║ 4} \\
\texttt{║◌◌◌│◌│O║ 5} \\
\texttt{╟─┬─┴─┼─╢ 6} \\
\texttt{║◌│◌◌◌│◌║ 7} \\
\texttt{╚═╧═══╧═╝} \\
\texttt{```} \\
\texttt{أنت تلعب حالياً الجولة 11 من أصل 28 كحد أقصى.} \\
\texttt{لقد جمعتما معاً 12 من أصل 16 عقوبة.} \\
\texttt{حرّك اللاعب الآخر كائناً على شبكته.} \\
\\ 
\\ 
\texttt{ما هو أمرك التالي؟} \\
            }
        }
    }
    & \\ \\

    \theutterance \stepcounter{utterance}  
    & & & \multicolumn{4}{p{0.6\linewidth}}{
        \cellcolor[rgb]{0.9,0.9,0.9}{
            \makecell[{{p{\linewidth}}}]{
                \texttt{\tiny{[GM$\langle$P2]}}
                \texttt{قل: على لوحتاي، الصف 1 امتلأ: P(1,1) S(2,1) A(3,1) M(4,1) D(6,1) T(7,1). بقي O في (7,5). الهدف O(5,1) لكن (5,1) مشغول بـ M. يمكن تحريك O إلى (5,1) بعد تحريك M؟ لكن M في مكانه المستهدف (5,1). على لوحة الآخر، الصف 1 لديه فراغات (6,1) و(7,1) فقط. اتفقنا على صف 3 كصف احتياطي. لنحرك O إلى (5,3) على لوحتاي؟ (5,3) كان A وتحرك إلى (4,1)، الآن (5,3) فارغ. لكن الهدف النهائي يجب أن يتطابق مع لوحة الآخر. اقترح: Nنضع O في (5,3) على كلا اللوحتين؟ أو نستخدم (4,3) لـ O. دعنا ننسق: هدف نهائي: P(1,1) S(2,1) A(3,1) M(4,1) D(5,1) T(6,1) O(7,1) على كلا اللوحتين. لكن على لوحتاي D في (6,1) و T في (7,1). يمكن تحريك D إلى (5,1) بعد إخلاء M؟ M في (4,1) الهدف M(4,1) جيد. D الهدف (5,1). T الهدف (6,1). O الهدف (7,1). إذن نحتاج إخلاء (5,1) و (6,1). حالياً M في (4,1) A في (3,1) S في (2,1) P في (1,1). (5,1) فارغ الآن؟ بعد تحريك M إلى (5,1) ثم إلى (4,1)؟ لا، M تحرك من (3,4) إلى (5,1) ثم؟ انتظر: M تحرك إلى (5,1) في الجولة 10. ثم لم يتحرك منذ ذلك. في اللوحة الحالية M في (5,1)? انظر اللوحة المحدثة: "P│SAMDT" يعني عمود1 P، 2 S، 3 A، 4 M، 5 D، 6 T؟ لا، "SAMDT" خمسة حروف في أعمدة 2{-}6؟ التنسيق: بعد "│" هناك "S◌D◌◌" أصلاً. الآن "SAMDT". هذا يعني عمود2 S، عمود3 A، عمود4 M، عمود5 D، عمود6 T. عمود7 فارغ. إذن M في (4,1)، D في (5,1)، T في (6,1). (7,1) فارغ. O في (7,5). هدف O(7,1). يمكن تحريك O إلى (7,1) مباشرة! جرب.} \\
            }
        }
    }
     \\ \\

    \theutterance \stepcounter{utterance}  
    & & & \multicolumn{2}{p{0.3\linewidth}}{
        \cellcolor[rgb]{0.9,0.9,0.9}{
            \makecell[{{p{\linewidth}}}]{
                \texttt{\tiny{[GM$|$GM]}}
                \texttt{قل: على لوحتاي، الصف 1 امتلأ: P(1,1) S(2,1) A(3,1) M(4,1) D(6,1) T(7,1). بقي O في (7,5). الهدف O(5,1) لكن (5,1) مشغول بـ M. يمكن تحريك O إلى (5,1) بعد تحريك M؟ لكن M في مكانه المستهدف (5,1). على لوحة الآخر، الصف 1 لديه فراغات (6,1) و(7,1) فقط. اتفقنا على صف 3 كصف احتياطي. لنحرك O إلى (5,3) على لوحتاي؟ (5,3) كان A وتحرك إلى (4,1)، الآن (5,3) فارغ. لكن الهدف النهائي يجب أن يتطابق مع لوحة الآخر. اقترح: Nنضع O في (5,3) على كلا اللوحتين؟ أو نستخدم (4,3) لـ O. دعنا ننسق: هدف نهائي: P(1,1) S(2,1) A(3,1) M(4,1) D(5,1) T(6,1) O(7,1) على كلا اللوحتين. لكن على لوحتاي D في (6,1) و T في (7,1). يمكن تحريك D إلى (5,1) بعد إخلاء M؟ M في (4,1) الهدف M(4,1) جيد. D الهدف (5,1). T الهدف (6,1). O الهدف (7,1). إذن نحتاج إخلاء (5,1) و (6,1). حالياً M في (4,1) A في (3,1) S في (2,1) P في (1,1). (5,1) فارغ الآن؟ بعد تحريك M إلى (5,1) ثم إلى (4,1)؟ لا، M تحرك من (3,4) إلى (5,1) ثم؟ انتظر: M تحرك إلى (5,1) في الجولة 10. ثم لم يتحرك منذ ذلك. في اللوحة الحالية M في (5,1)? انظر اللوحة المحدثة: "P│SAMDT" يعني عمود1 P، 2 S، 3 A، 4 M، 5 D، 6 T؟ لا، "SAMDT" خمسة حروف في أعمدة 2{-}6؟ التنسيق: بعد "│" هناك "S◌D◌◌" أصلاً. الآن "SAMDT". هذا يعني عمود2 S، عمود3 A، عمود4 M، عمود5 D، عمود6 T. عمود7 فارغ. إذن M في (4,1)، D في (5,1)، T في (6,1). (7,1) فارغ. O في (7,5). هدف O(7,1). يمكن تحريك O إلى (7,1) مباشرة! جرب.} \\
            }
        }
    }
    & & \\ \\

    \theutterance \stepcounter{utterance}  
    & & \multicolumn{4}{p{0.6\linewidth}}{
        \cellcolor[rgb]{0.9,0.9,0.9}{
            \makecell[{{p{\linewidth}}}]{
                \texttt{\tiny{[P1$\langle$GM]}}
                \texttt{تم تحريك 'T' إلى (6,1) بنجاح. تبدو شبكتك المحدّثة هكذا:} \\
\texttt{```} \\
\texttt{ 1234567} \\
\texttt{╔═╤═════╗} \\
\texttt{║P│OAST◌║ 1} \\
\texttt{╟─┤◌◌◌◌◌║ 2} \\
\texttt{║◌└─┐◌◌◌║ 3} \\
\texttt{║◌◌◌├─┐◌║ 4} \\
\texttt{║◌◌◌│◌│◌║ 5} \\
\texttt{╟─┬─┴─┼─╢ 6} \\
\texttt{║◌│◌DM│◌║ 7} \\
\texttt{╚═╧═══╧═╝} \\
\texttt{```} \\
\texttt{أنت تلعب حالياً الجولة 11 من أصل 28 كحد أقصى.} \\
\texttt{لقد جمعتما معاً 12 من أصل 16 عقوبة.} \\
\texttt{أرسل اللاعب الآخر هذه الرسالة:} \\
\\ 
\\ 
\texttt{"على لوحتاي، الصف 1 امتلأ: P(1,1) S(2,1) A(3,1) M(4,1) D(6,1) T(7,1). بقي O في (7,5). الهدف O(5,1) لكن (5,1) مشغول بـ M. يمكن تحريك O إلى (5,1) بعد تحريك M؟ لكن M في مكانه المستهدف (5,1). على لوحة الآخر، الصف 1 لديه فراغات (6,1) و(7,1) فقط. اتفقنا على صف 3 كصف احتياطي. لنحرك O إلى (5,3) على لوحتاي؟ (5,3) كان A وتحرك إلى (4,1)، الآن (5,3) فارغ. لكن الهدف النهائي يجب أن يتطابق مع لوحة الآخر. اقترح: Nنضع O في (5,3) على كلا اللوحتين؟ أو نستخدم (4,3) لـ O. دعنا ننسق: هدف نهائي: P(1,1) S(2,1) A(3,1) M(4,1) D(5,1) T(6,1) O(7,1) على كلا اللوحتين. لكن على لوحتاي D في (6,1) و T في (7,1). يمكن تحريك D إلى (5,1) بعد إخلاء M؟ M في (4,1) الهدف M(4,1) جيد. D الهدف (5,1). T الهدف (6,1). O الهدف (7,1). إذن نحتاج إخلاء (5,1) و (6,1). حالياً M في (4,1) A في (3,1) S في (2,1) P في (1,1). (5,1) فارغ الآن؟ بعد تحريك M إلى (5,1) ثم إلى (4,1)؟ لا، M تحرك من (3,4) إلى (5,1) ثم؟ انتظر: M تحرك إلى (5,1) في الجولة 10. ثم لم يتحرك منذ ذلك. في اللوحة الحالية M في (5,1)? انظر اللوحة المحدثة: "P│SAMDT" يعني عمود1 P، 2 S، 3 A، 4 M، 5 D، 6 T؟ لا، "SAMDT" خمسة حروف في أعمدة 2{-}6؟ التنسيق: بعد "│" هناك "S◌D◌◌" أصلاً. الآن "SAMDT". هذا يعني عمود2 S، عمود3 A، عمود4 M، عمود5 D، عمود6 T. عمود7 فارغ. إذن M في (4,1)، D في (5,1)، T في (6,1). (7,1) فارغ. O في (7,5). هدف O(7,1). يمكن تحريك O إلى (7,1) مباشرة! جرب."} \\
\\ 
\\ 
\texttt{ما هو أمرك التالي؟} \\
            }
        }
    }
    & \\ \\

    \theutterance \stepcounter{utterance}  
    & \multicolumn{4}{p{0.6\linewidth}}{
        \cellcolor[rgb]{0.9,0.9,0.9}{
            \makecell[{{p{\linewidth}}}]{
                \texttt{\tiny{[P1$\rangle$GM]}}
                \texttt{قل: فهمت خطة اللاعب 2. على لوحتي، الصف 1 الآن: P(1,1) O(3,1) A(3,1؟ لا A في (3,1) أصلاً؟ انظر اللوحة: "P│OAST◌" يعني عمود1 P، عمود2 جدار، عمود3 O، عمود4 A، عمود5 S، عمود6 T، عمود7 فارغ. إذن A في (4,1)، S في (5,1)، T في (6,1). (7,1) فارغ. الكائنات المتبقية على لوحتي: D في (5,7)، M في (6,7). سأحرك A إلى (3,1) لكن O هناك. يجب تحريك O بعيداً. الأفضل: نستخدم الصف 3 للكائنات المتبقية. هدف نهائي مقترح لكلا اللوحتين:} \\
\texttt{P(1,1) S(2,1) O(3,1) A(4,1) T(5,1) D(6,1) M(7,1) {-} لكن (2,1) جدار على لوحتي.} \\
\texttt{بديل: نستخدم إحداثيات مختلفة تماماً. أقترح الهدف: P(1,1) O(3,1) S(5,1) T(6,1) A(4,3) D(5,3) M(6,3) {-} ولكن يجب أن يتطابق مع لاعب 2.} \\
\texttt{لنسق بدقة: أخبرني بإحداثيات الهدف النهائية التي وصلت لها على لوحتك، وسأطابقها.} \\
            }
        }
    }
    & & \\ \\

    \theutterance \stepcounter{utterance}  
    & & & \multicolumn{2}{p{0.3\linewidth}}{
        \cellcolor[rgb]{0.9,0.9,0.9}{
            \makecell[{{p{\linewidth}}}]{
                \texttt{\tiny{[GM$|$GM]}}
                \texttt{قل: فهمت خطة اللاعب 2. على لوحتي، الصف 1 الآن: P(1,1) O(3,1) A(3,1؟ لا A في (3,1) أصلاً؟ انظر اللوحة: "P│OAST◌" يعني عمود1 P، عمود2 جدار، عمود3 O، عمود4 A، عمود5 S، عمود6 T، عمود7 فارغ. إذن A في (4,1)، S في (5,1)، T في (6,1). (7,1) فارغ. الكائنات المتبقية على لوحتي: D في (5,7)، M في (6,7). سأحرك A إلى (3,1) لكن O هناك. يجب تحريك O بعيداً. الأفضل: نستخدم الصف 3 للكائنات المتبقية. هدف نهائي مقترح لكلا اللوحتين:} \\
\texttt{P(1,1) S(2,1) O(3,1) A(4,1) T(5,1) D(6,1) M(7,1) {-} لكن (2,1) جدار على لوحتي.} \\
\texttt{بديل: نستخدم إحداثيات مختلفة تماماً. أقترح الهدف: P(1,1) O(3,1) S(5,1) T(6,1) A(4,3) D(5,3) M(6,3) {-} ولكن يجب أن يتطابق مع لاعب 2.} \\
\texttt{لنسق بدقة: أخبرني بإحداثيات الهدف النهائية التي وصلت لها على لوحتك، وسأطابقها.} \\
            }
        }
    }
    & & \\ \\

    \theutterance \stepcounter{utterance}  
    & & & \multicolumn{2}{p{0.3\linewidth}}{
        \cellcolor[rgb]{0.9,0.9,0.9}{
            \makecell[{{p{\linewidth}}}]{
                \texttt{\tiny{[GM$|$GM]}}
                \texttt{Invalid format: tail is not empty: '} \\
\texttt{P(1,1) S(2,1) O(3,1) A(4,1) T(5,1) D(6,1) M(7,1) {-} لكن (2,1) جدار على لوحتي.} \\
\texttt{بديل: نستخدم إحداثيات مختلفة تماماً. أقترح الهدف: P(1,1) O(3,1) S(5,1) T(6,1) A(4,3) D(5,3) M(6,3) {-} ولكن يجب أن يتطابق مع لاعب 2.} \\
\texttt{لنسق بدقة: أخبرني بإحداثيات الهدف النهائية التي وصلت لها على لوحتك، وسأطابقها.} \\
            }
        }
    }
    & & \\ \\

    \theutterance \stepcounter{utterance}  
    & & \multicolumn{4}{p{0.6\linewidth}}{
        \cellcolor[rgb]{0.9,0.9,0.9}{
            \makecell[{{p{\linewidth}}}]{
                \texttt{\tiny{[P1$\langle$GM]}}
                \texttt{عقوبة: يجب ألا تحتوي رسالتك على أي شيء بعد الأمر!} \\
\\ 
\texttt{تأكد من أن ردك يحتوي فقط على `قل: <رسالة>` أو `حرّك: <كائن>, (<X>, <Y>)`، ولا شيء غير ذلك!} \\
\\ 
\texttt{لقد جمعتما معاً 13 من أصل 16 عقوبة.} \\
\texttt{يرجى المحاولة مرة أخرى!} \\
            }
        }
    }
    & \\ \\

    \theutterance \stepcounter{utterance}  
    & \multicolumn{4}{p{0.6\linewidth}}{
        \cellcolor[rgb]{0.9,0.9,0.9}{
            \makecell[{{p{\linewidth}}}]{
                \texttt{\tiny{[P1$\rangle$GM]}}
                \texttt{قل: فهمت خطة اللاعب 2. على لوحتي، الصف 1: P(1,1) O(3,1) A(4,1) S(5,1) T(6,1) (7,1) فارغ. المتبقي: D(5,7) M(6,7). هدف نهائي مقترح لكلا اللوحتين: P(1,1) S(2,1) O(3,1) A(4,1) T(5,1) D(6,1) M(7,1) {-} لكن (2,1) جدار على لوحتي. لنسق: أخبرني بإحداثيات هدفك النهائية لأطابقها.} \\
            }
        }
    }
    & & \\ \\

    \theutterance \stepcounter{utterance}  
    & & & \multicolumn{2}{p{0.3\linewidth}}{
        \cellcolor[rgb]{0.9,0.9,0.9}{
            \makecell[{{p{\linewidth}}}]{
                \texttt{\tiny{[GM$|$GM]}}
                \texttt{قل: فهمت خطة اللاعب 2. على لوحتي، الصف 1: P(1,1) O(3,1) A(4,1) S(5,1) T(6,1) (7,1) فارغ. المتبقي: D(5,7) M(6,7). هدف نهائي مقترح لكلا اللوحتين: P(1,1) S(2,1) O(3,1) A(4,1) T(5,1) D(6,1) M(7,1) {-} لكن (2,1) جدار على لوحتي. لنسق: أخبرني بإحداثيات هدفك النهائية لأطابقها.} \\
            }
        }
    }
    & & \\ \\

    \theutterance \stepcounter{utterance}  
    & & \multicolumn{4}{p{0.6\linewidth}}{
        \cellcolor[rgb]{0.9,0.9,0.9}{
            \makecell[{{p{\linewidth}}}]{
                \texttt{\tiny{[GM$\rangle$P2]}}
                \texttt{تم تمرير رسالتك إلى اللاعب الآخر.} \\
\\ 
\\ 
\texttt{أنت تلعب حالياً الجولة 12 من أصل 28 كحد أقصى.} \\
\texttt{لقد جمعتما معاً 13 من أصل 16 عقوبة.} \\
\texttt{أرسل اللاعب الآخر هذه الرسالة:} \\
\\ 
\\ 
\texttt{"فهمت خطة اللاعب 2. على لوحتي، الصف 1: P(1,1) O(3,1) A(4,1) S(5,1) T(6,1) (7,1) فارغ. المتبقي: D(5,7) M(6,7). هدف نهائي مقترح لكلا اللوحتين: P(1,1) S(2,1) O(3,1) A(4,1) T(5,1) D(6,1) M(7,1) {-} لكن (2,1) جدار على لوحتي. لنسق: أخبرني بإحداثيات هدفك النهائية لأطابقها."} \\
\\ 
\\ 
\texttt{ما هو أمرك التالي؟} \\
            }
        }
    }
    & \\ \\

    \theutterance \stepcounter{utterance}  
    & & & \multicolumn{4}{p{0.6\linewidth}}{
        \cellcolor[rgb]{0.9,0.9,0.9}{
            \makecell[{{p{\linewidth}}}]{
                \texttt{\tiny{[GM$\langle$P2]}}
                \texttt{حرّك: O, (7, 1)} \\
            }
        }
    }
     \\ \\

    \theutterance \stepcounter{utterance}  
    & & & \multicolumn{2}{p{0.3\linewidth}}{
        \cellcolor[rgb]{0.9,0.9,0.9}{
            \makecell[{{p{\linewidth}}}]{
                \texttt{\tiny{[GM$|$GM]}}
                \texttt{حرّك: O, (7, 1)} \\
            }
        }
    }
    & & \\ \\

    \theutterance \stepcounter{utterance}  
    & & & \multicolumn{2}{p{0.3\linewidth}}{
        \cellcolor[rgb]{0.9,0.9,0.9}{
            \makecell[{{p{\linewidth}}}]{
                \texttt{\tiny{[GM$|$GM]}}
                \texttt{عقوبة: (7,1) ليست فارغة، بل تحتوي على 'T'.} \\
\\ 
\texttt{لقد جمعتما معاً 14 من أصل 16 عقوبة.} \\
\texttt{يرجى المحاولة مرة أخرى!} \\
            }
        }
    }
    & & \\ \\

    \theutterance \stepcounter{utterance}  
    & & \multicolumn{4}{p{0.6\linewidth}}{
        \cellcolor[rgb]{0.9,0.9,0.9}{
            \makecell[{{p{\linewidth}}}]{
                \texttt{\tiny{[GM$\rangle$P2]}}
                \texttt{عقوبة: (7,1) ليست فارغة، بل تحتوي على 'T'.} \\
\\ 
\texttt{لقد جمعتما معاً 14 من أصل 16 عقوبة.} \\
\texttt{يرجى المحاولة مرة أخرى!} \\
            }
        }
    }
    & \\ \\

    \theutterance \stepcounter{utterance}  
    & & & \multicolumn{4}{p{0.6\linewidth}}{
        \cellcolor[rgb]{0.9,0.9,0.9}{
            \makecell[{{p{\linewidth}}}]{
                \texttt{\tiny{[GM$\langle$P2]}}
                \texttt{قل: على لوحتاي T في (6,1) و (7,1) فارغ. لكن النظام يقول (7,1) يحتوي على T. ربما تحديث اللوحة: "P│SAMDT" يعني الأعمدة 2{-}6: S A M D T. إذن T في العمود 6. العمود 7 فارغ. الإحداثيات (7,1) يجب أن تكون فارغة. جرب مرة أخرى.} \\
            }
        }
    }
     \\ \\

    \theutterance \stepcounter{utterance}  
    & & & \multicolumn{2}{p{0.3\linewidth}}{
        \cellcolor[rgb]{0.9,0.9,0.9}{
            \makecell[{{p{\linewidth}}}]{
                \texttt{\tiny{[GM$|$GM]}}
                \texttt{قل: على لوحتاي T في (6,1) و (7,1) فارغ. لكن النظام يقول (7,1) يحتوي على T. ربما تحديث اللوحة: "P│SAMDT" يعني الأعمدة 2{-}6: S A M D T. إذن T في العمود 6. العمود 7 فارغ. الإحداثيات (7,1) يجب أن تكون فارغة. جرب مرة أخرى.} \\
            }
        }
    }
    & & \\ \\

    \theutterance \stepcounter{utterance}  
    & & \multicolumn{4}{p{0.6\linewidth}}{
        \cellcolor[rgb]{0.9,0.9,0.9}{
            \makecell[{{p{\linewidth}}}]{
                \texttt{\tiny{[P1$\langle$GM]}}
                \texttt{تم تمرير رسالتك إلى اللاعب الآخر.} \\
\\ 
\\ 
\texttt{أنت تلعب حالياً الجولة 12 من أصل 28 كحد أقصى.} \\
\texttt{لقد جمعتما معاً 14 من أصل 16 عقوبة.} \\
\texttt{أرسل اللاعب الآخر هذه الرسالة:} \\
\\ 
\\ 
\texttt{"على لوحتاي T في (6,1) و (7,1) فارغ. لكن النظام يقول (7,1) يحتوي على T. ربما تحديث اللوحة: "P│SAMDT" يعني الأعمدة 2{-}6: S A M D T. إذن T في العمود 6. العمود 7 فارغ. الإحداثيات (7,1) يجب أن تكون فارغة. جرب مرة أخرى."} \\
\\ 
\\ 
\texttt{ما هو أمرك التالي؟} \\
            }
        }
    }
    & \\ \\

    \theutterance \stepcounter{utterance}  
    & \multicolumn{4}{p{0.6\linewidth}}{
        \cellcolor[rgb]{0.9,0.9,0.9}{
            \makecell[{{p{\linewidth}}}]{
                \texttt{\tiny{[P1$\rangle$GM]}}
                \texttt{قل: على لوحتك T في (6,1) و(7,1) فارغ. على لوحتي (7,1) فارغ. س أحرك D إلى (7,1) على لوحتي، ثم M إلى خلية في الصف 3. لكننا بحاجة لتطابق تام. ما هي إحداثياتك النهائية لكل كائن الآن؟ P(1,1) S(2,1) A(3,1) M(4,1) D(5,1) T(6,1) O(7,1)؟} \\
            }
        }
    }
    & & \\ \\

    \theutterance \stepcounter{utterance}  
    & & & \multicolumn{2}{p{0.3\linewidth}}{
        \cellcolor[rgb]{0.9,0.9,0.9}{
            \makecell[{{p{\linewidth}}}]{
                \texttt{\tiny{[GM$|$GM]}}
                \texttt{قل: على لوحتك T في (6,1) و(7,1) فارغ. على لوحتي (7,1) فارغ. س أحرك D إلى (7,1) على لوحتي، ثم M إلى خلية في الصف 3. لكننا بحاجة لتطابق تام. ما هي إحداثياتك النهائية لكل كائن الآن؟ P(1,1) S(2,1) A(3,1) M(4,1) D(5,1) T(6,1) O(7,1)؟} \\
            }
        }
    }
    & & \\ \\

    \theutterance \stepcounter{utterance}  
    & & \multicolumn{4}{p{0.6\linewidth}}{
        \cellcolor[rgb]{0.9,0.9,0.9}{
            \makecell[{{p{\linewidth}}}]{
                \texttt{\tiny{[GM$\rangle$P2]}}
                \texttt{تم تمرير رسالتك إلى اللاعب الآخر.} \\
\\ 
\\ 
\texttt{أنت تلعب حالياً الجولة 13 من أصل 28 كحد أقصى.} \\
\texttt{لقد جمعتما معاً 14 من أصل 16 عقوبة.} \\
\texttt{أرسل اللاعب الآخر هذه الرسالة:} \\
\\ 
\\ 
\texttt{"على لوحتك T في (6,1) و(7,1) فارغ. على لوحتي (7,1) فارغ. س أحرك D إلى (7,1) على لوحتي، ثم M إلى خلية في الصف 3. لكننا بحاجة لتطابق تام. ما هي إحداثياتك النهائية لكل كائن الآن؟ P(1,1) S(2,1) A(3,1) M(4,1) D(5,1) T(6,1) O(7,1)؟"} \\
\\ 
\\ 
\texttt{ما هو أمرك التالي؟} \\
            }
        }
    }
    & \\ \\

    \theutterance \stepcounter{utterance}  
    & & & \multicolumn{4}{p{0.6\linewidth}}{
        \cellcolor[rgb]{0.9,0.9,0.9}{
            \makecell[{{p{\linewidth}}}]{
                \texttt{\tiny{[GM$\langle$P2]}}
                \texttt{حرّك: O, (7, 1)} \\
            }
        }
    }
     \\ \\

    \theutterance \stepcounter{utterance}  
    & & & \multicolumn{2}{p{0.3\linewidth}}{
        \cellcolor[rgb]{0.9,0.9,0.9}{
            \makecell[{{p{\linewidth}}}]{
                \texttt{\tiny{[GM$|$GM]}}
                \texttt{حرّك: O, (7, 1)} \\
            }
        }
    }
    & & \\ \\

    \theutterance \stepcounter{utterance}  
    & & & \multicolumn{2}{p{0.3\linewidth}}{
        \cellcolor[rgb]{0.9,0.9,0.9}{
            \makecell[{{p{\linewidth}}}]{
                \texttt{\tiny{[GM$|$GM]}}
                \texttt{عقوبة: (7,1) ليست فارغة، بل تحتوي على 'T'.} \\
\\ 
\texttt{لقد جمعتما معاً 15 من أصل 16 عقوبة.} \\
\texttt{يرجى المحاولة مرة أخرى!} \\
            }
        }
    }
    & & \\ \\

    \theutterance \stepcounter{utterance}  
    & & \multicolumn{4}{p{0.6\linewidth}}{
        \cellcolor[rgb]{0.9,0.9,0.9}{
            \makecell[{{p{\linewidth}}}]{
                \texttt{\tiny{[GM$\rangle$P2]}}
                \texttt{عقوبة: (7,1) ليست فارغة، بل تحتوي على 'T'.} \\
\\ 
\texttt{لقد جمعتما معاً 15 من أصل 16 عقوبة.} \\
\texttt{يرجى المحاولة مرة أخرى!} \\
            }
        }
    }
    & \\ \\

    \theutterance \stepcounter{utterance}  
    & & & \multicolumn{4}{p{0.6\linewidth}}{
        \cellcolor[rgb]{0.9,0.9,0.9}{
            \makecell[{{p{\linewidth}}}]{
                \texttt{\tiny{[GM$\langle$P2]}}
                \texttt{قل: نظامي يظهر T في (6,1). اللوحة: "P│SAMDT" {-} هذا 5 حروف بعد الفاصل. الفاصل في عمود1؟ تنسيق اللوحة: عمود1 P، ثم فاصل "│"، ثم عمود2 S، عمود3 A، عمود4 M، عمود5 D، عمود6 T، عمود7 فارغ. إذن T في (6,1). (7,1) فارغ. لماذا الخطأ؟ ربما الفهرسة: الإحداثيات (س، ص) حيث س العمود، ص الصف. الصف 1. جرب تحريك O إلى (6,1)؟ لا، الهدف (7,1). جرب تحريك T إلى (7,1) أولاً لإخلاء (6,1) لـ D؟ لكن D في (5,1) الهدف D(5,1) جيد. T الهدف (6,1) جيد. O الهدف (7,1). فقط O يحتاج حركة. جرب مرة أخرى حرّك: O, (7, 1) بعناية.} \\
            }
        }
    }
     \\ \\

    \theutterance \stepcounter{utterance}  
    & & & \multicolumn{2}{p{0.3\linewidth}}{
        \cellcolor[rgb]{0.9,0.9,0.9}{
            \makecell[{{p{\linewidth}}}]{
                \texttt{\tiny{[GM$|$GM]}}
                \texttt{قل: نظامي يظهر T في (6,1). اللوحة: "P│SAMDT" {-} هذا 5 حروف بعد الفاصل. الفاصل في عمود1؟ تنسيق اللوحة: عمود1 P، ثم فاصل "│"، ثم عمود2 S، عمود3 A، عمود4 M، عمود5 D، عمود6 T، عمود7 فارغ. إذن T في (6,1). (7,1) فارغ. لماذا الخطأ؟ ربما الفهرسة: الإحداثيات (س، ص) حيث س العمود، ص الصف. الصف 1. جرب تحريك O إلى (6,1)؟ لا، الهدف (7,1). جرب تحريك T إلى (7,1) أولاً لإخلاء (6,1) لـ D؟ لكن D في (5,1) الهدف D(5,1) جيد. T الهدف (6,1) جيد. O الهدف (7,1). فقط O يحتاج حركة. جرب مرة أخرى حرّك: O, (7, 1) بعناية.} \\
            }
        }
    }
    & & \\ \\

    \theutterance \stepcounter{utterance}  
    & & & \multicolumn{2}{p{0.3\linewidth}}{
        \cellcolor[rgb]{0.9,0.9,0.9}{
            \makecell[{{p{\linewidth}}}]{
                \texttt{\tiny{[GM$|$GM]}}
                \texttt{Invalid response format: قل: نظامي يظهر T في (6,1). اللوحة: "P│SAMDT" {-} هذا 5 حروف بعد الفاصل. الفاصل في عمود1؟ تنسيق اللوحة: عمود1 P، ثم فاصل "│"، ثم عمود2 S، عمود3 A، عمود4 M، عمود5 D، عمود6 T، عمود7 فارغ. إذن T في (6,1). (7,1) فارغ. لماذا الخطأ؟ ربما الفهرسة: الإحداثيات (س، ص) حيث س العمود، ص الصف. الصف 1. جرب تحريك O إلى (6,1)؟ لا، الهدف (7,1). جرب تحريك T إلى (7,1) أولاً لإخلاء (6,1) لـ D؟ لكن D في (5,1) الهدف D(5,1) جيد. T الهدف (6,1) جيد. O الهدف (7,1). فقط O يحتاج حركة. جرب مرة أخرى حرّك: O, (7, 1) بعناية.} \\
            }
        }
    }
    & & \\ \\

    \theutterance \stepcounter{utterance}  
    & & \multicolumn{4}{p{0.6\linewidth}}{
        \cellcolor[rgb]{0.9,0.9,0.9}{
            \makecell[{{p{\linewidth}}}]{
                \texttt{\tiny{[GM$\rangle$P2]}}
                \texttt{عقوبة: تحتوي رسالتك على أكثر من أمر واحد!} \\
\\ 
\texttt{تأكد من أن ردك يحتوي فقط على `قل: <رسالة>` أو `حرّك: <كائن>, (<X>, <Y>)`، ولا شيء غير ذلك!} \\
\\ 
\texttt{لقد جمعتما معاً 16 من أصل 16 عقوبة.} \\
\texttt{يرجى المحاولة مرة أخرى!} \\
            }
        }
    }
    & \\ \\

    \theutterance \stepcounter{utterance}  
    & & & \multicolumn{4}{p{0.6\linewidth}}{
        \cellcolor[rgb]{0.9,0.9,0.9}{
            \makecell[{{p{\linewidth}}}]{
                \texttt{\tiny{[GM$\langle$P2]}}
                \texttt{قل: نظامي يظهر T في (6,1). اللوحة: P│SAMDT {-} هذا 5 حروف بعد الفاصل. الفاصل في عمود1؟ تنسيق اللوحة: عمود1 P، ثم فاصل │، ثم عمود2 S، عمود3 A، عمود4 M، عمود5 D، عمود6 T، عمود7 فارغ. إذن T في (6,1). (7,1) فارغ. لماذا الخطأ؟} \\
            }
        }
    }
     \\ \\

    \theutterance \stepcounter{utterance}  
    & & & \multicolumn{2}{p{0.3\linewidth}}{
        \cellcolor[rgb]{0.9,0.9,0.9}{
            \makecell[{{p{\linewidth}}}]{
                \texttt{\tiny{[GM$|$GM]}}
                \texttt{قل: نظامي يظهر T في (6,1). اللوحة: P│SAMDT {-} هذا 5 حروف بعد الفاصل. الفاصل في عمود1؟ تنسيق اللوحة: عمود1 P، ثم فاصل │، ثم عمود2 S، عمود3 A، عمود4 M، عمود5 D، عمود6 T، عمود7 فارغ. إذن T في (6,1). (7,1) فارغ. لماذا الخطأ؟} \\
            }
        }
    }
    & & \\ \\

    \theutterance \stepcounter{utterance}  
    & & \multicolumn{4}{p{0.6\linewidth}}{
        \cellcolor[rgb]{0.9,0.9,0.9}{
            \makecell[{{p{\linewidth}}}]{
                \texttt{\tiny{[P1$\langle$GM]}}
                \texttt{تم تمرير رسالتك إلى اللاعب الآخر.} \\
\\ 
\\ 
\texttt{أنت تلعب حالياً الجولة 13 من أصل 28 كحد أقصى.} \\
\texttt{لقد جمعتما معاً 16 من أصل 16 عقوبة.} \\
\texttt{أرسل اللاعب الآخر هذه الرسالة:} \\
\\ 
\\ 
\texttt{"نظامي يظهر T في (6,1). اللوحة: P│SAMDT {-} هذا 5 حروف بعد الفاصل. الفاصل في عمود1؟ تنسيق اللوحة: عمود1 P، ثم فاصل │، ثم عمود2 S، عمود3 A، عمود4 M، عمود5 D، عمود6 T، عمود7 فارغ. إذن T في (6,1). (7,1) فارغ. لماذا الخطأ؟"} \\
\\ 
\\ 
\texttt{ما هو أمرك التالي؟} \\
            }
        }
    }
    & \\ \\

    \theutterance \stepcounter{utterance}  
    & \multicolumn{4}{p{0.6\linewidth}}{
        \cellcolor[rgb]{0.9,0.9,0.9}{
            \makecell[{{p{\linewidth}}}]{
                \texttt{\tiny{[P1$\rangle$GM]}}
                \texttt{لقد وصلنا إلى 16 عقوبة (الحد الأقصى). اللعبة انتهت بخسارة.} \\
\\ 
\texttt{مع ذلك، سأحاول مرة أخيرة لإنهاء اللعبة بنجاح إن أمكن:} \\
\\ 
\texttt{قل: على لوحتك T في (6,1) و(7,1) فارغ. على لوحتي (7,1) فارغ. س أحرك D إلى (7,1) على لوحتي. ثم نحتاج لمطابقة M و O. على لوحتك O في (7,5) الهدف (7,1). على لوحتي O في (3,1). هل يمكننا إنهاء اللعبة بمواقع: P(1,1) S(2,1؟ جدار) O(3,1) A(4,1) T(5,1؟ لا T في 6,1) D(7,1) M(؟) {-} مستحيل مع الجدار في (2,1). نحتاج خطة جديدة.} \\
            }
        }
    }
    & & \\ \\

    \theutterance \stepcounter{utterance}  
    & & & \multicolumn{2}{p{0.3\linewidth}}{
        \cellcolor[rgb]{0.9,0.9,0.9}{
            \makecell[{{p{\linewidth}}}]{
                \texttt{\tiny{[GM$|$GM]}}
                \texttt{لقد وصلنا إلى 16 عقوبة (الحد الأقصى). اللعبة انتهت بخسارة.} \\
\\ 
\texttt{مع ذلك، سأحاول مرة أخيرة لإنهاء اللعبة بنجاح إن أمكن:} \\
\\ 
\texttt{قل: على لوحتك T في (6,1) و(7,1) فارغ. على لوحتي (7,1) فارغ. س أحرك D إلى (7,1) على لوحتي. ثم نحتاج لمطابقة M و O. على لوحتك O في (7,5) الهدف (7,1). على لوحتي O في (3,1). هل يمكننا إنهاء اللعبة بمواقع: P(1,1) S(2,1؟ جدار) O(3,1) A(4,1) T(5,1؟ لا T في 6,1) D(7,1) M(؟) {-} مستحيل مع الجدار في (2,1). نحتاج خطة جديدة.} \\
            }
        }
    }
    & & \\ \\

    \theutterance \stepcounter{utterance}  
    & & & \multicolumn{2}{p{0.3\linewidth}}{
        \cellcolor[rgb]{0.9,0.9,0.9}{
            \makecell[{{p{\linewidth}}}]{
                \texttt{\tiny{[GM$|$GM]}}
                \texttt{Invalid format: head is not empty: 'لقد وصلنا إلى 16 عقوبة (الحد الأقصى). اللعبة انتهت بخسارة.} \\
\\ 
\texttt{مع ذلك، سأحاول مرة أخيرة لإنهاء اللعبة بنجاح إن أمكن:} \\
\\ 
\texttt{'} \\
            }
        }
    }
    & & \\ \\

    \theutterance \stepcounter{utterance}  
    & & & \multicolumn{2}{p{0.3\linewidth}}{
        \cellcolor[rgb]{0.9,0.9,0.9}{
            \makecell[{{p{\linewidth}}}]{
                \texttt{\tiny{[GM$|$GM]}}
                \texttt{Maximum number of penalties exceeded} \\
            }
        }
    }
    & & \\ \\

    \theutterance \stepcounter{utterance}  
    & & & \multicolumn{2}{p{0.3\linewidth}}{
        \cellcolor[rgb]{0.9,0.9,0.9}{
            \makecell[{{p{\linewidth}}}]{
                \texttt{\tiny{[GM$|$GM]}}
                \texttt{Initial states:} \\
\texttt{```} \\
\texttt{Player 1:} \\
\texttt{ 1234567} \\
\texttt{╔═╤═════╗} \\
\texttt{║◌│◌A◌◌◌║ 1} \\
\texttt{╟─┤◌T◌◌◌║ 2} \\
\texttt{║O└─┐◌◌◌║ 3} \\
\texttt{║◌◌◌├─┐◌║ 4} \\
\texttt{║◌◌S│◌│◌║ 5} \\
\texttt{╟─┬─┴─┼─╢ 6} \\
\texttt{║P│◌DM│◌║ 7} \\
\texttt{╚═╧═══╧═╝} \\
\\ 
\\ 
\texttt{Player 2:} \\
\texttt{ 1234567} \\
\texttt{╔═╤═════╗} \\
\texttt{║P│S◌D◌◌║ 1} \\
\texttt{╟─┤◌◌◌◌◌║ 2} \\
\texttt{║◌└─┐◌A◌║ 3} \\
\texttt{║◌◌M├─┐◌║ 4} \\
\texttt{║◌T◌│◌│O║ 5} \\
\texttt{╟─┬─┴─┼─╢ 6} \\
\texttt{║◌│◌◌◌│◌║ 7} \\
\texttt{╚═╧═══╧═╝} \\
\\ 
\texttt{```} \\
            }
        }
    }
    & & \\ \\

    \theutterance \stepcounter{utterance}  
    & & & \multicolumn{2}{p{0.3\linewidth}}{
        \cellcolor[rgb]{0.9,0.9,0.9}{
            \makecell[{{p{\linewidth}}}]{
                \texttt{\tiny{[GM$|$GM]}}
                \texttt{End states:} \\
\texttt{```} \\
\texttt{Player 1:} \\
\texttt{ 1234567} \\
\texttt{╔═╤═════╗} \\
\texttt{║P│OAST◌║ 1} \\
\texttt{╟─┤◌◌◌◌◌║ 2} \\
\texttt{║◌└─┐◌◌◌║ 3} \\
\texttt{║◌◌◌├─┐◌║ 4} \\
\texttt{║◌◌◌│◌│◌║ 5} \\
\texttt{╟─┬─┴─┼─╢ 6} \\
\texttt{║◌│◌DM│◌║ 7} \\
\texttt{╚═╧═══╧═╝} \\
\\ 
\\ 
\texttt{Player 2:} \\
\texttt{ 1234567} \\
\texttt{╔═╤═════╗} \\
\texttt{║P│SAMDT║ 1} \\
\texttt{╟─┤◌◌◌◌◌║ 2} \\
\texttt{║◌└─┐◌◌◌║ 3} \\
\texttt{║◌◌◌├─┐◌║ 4} \\
\texttt{║◌◌◌│◌│O║ 5} \\
\texttt{╟─┬─┴─┼─╢ 6} \\
\texttt{║◌│◌◌◌│◌║ 7} \\
\texttt{╚═╧═══╧═╝} \\
\\ 
\texttt{```} \\
            }
        }
    }
    & & \\ \\

    \theutterance \stepcounter{utterance}  
    & & & \multicolumn{2}{p{0.3\linewidth}}{
        \cellcolor[rgb]{0.9,0.9,0.9}{
            \makecell[{{p{\linewidth}}}]{
                \texttt{\tiny{[GM$|$GM]}}
                \texttt{* success: False} \\
\texttt{* lose: False} \\
\texttt{* aborted: True} \\
\texttt{{-}{-}{-}{-}{-}{-}{-}} \\
\texttt{* Shifts: 6.00} \\
\texttt{* Max Shifts: 12.00} \\
\texttt{* Min Shifts: 6.00} \\
\texttt{* End Distance Sum: 20.98} \\
\texttt{* Init Distance Sum: 32.45} \\
\texttt{* Expected Distance Sum: 29.33} \\
\texttt{* Penalties: 17.00} \\
\texttt{* Max Penalties: 16.00} \\
\texttt{* Rounds: 13.00} \\
\texttt{* Max Rounds: 28.00} \\
\texttt{* Object Count: 7.00} \\
            }
        }
    }
    & & \\ \\

\end{supertabular}
}

\end{document}
